% sage_latex_guidelines.tex V1.20, 14 January 2017

\documentclass[Afour,sagev,times]{sagej}

\usepackage{moreverb,url}
\usepackage{amsmath}
\usepackage{graphicx}
\usepackage{rotating}
\usepackage[table]{xcolor}
\usepackage{booktabs}
\usepackage{array}
\usepackage{comment}
\usepackage[T1]{fontenc}
\usepackage[utf8]{inputenc}
\usepackage[colorlinks,bookmarksopen,bookmarksnumbered,citecolor=red,urlcolor=red]{hyperref}

\newcommand\BibTeX{{\rmfamily B\kern-.05em \textsc{i\kern-.025em b}\kern-.08em
T\kern-.1667em\lower.7ex\hbox{E}\kern-.125emX}}

\newcommand{\abbrevnote}[1]{\par\vspace{0.35em}{\footnotesize\textbf{Abbreviations:} #1\par}}

\def\volumeyear{2016}

\begin{document}

\runninghead{Computational Framework for nsPEF-derived EV Engineering}

\title{A Modular Systems-Biology Framework for nsPEF-Enabled Extracellular Vesicle Engineering}

\author{Arash Ghorbannia\affilnum{1,2}, Mihir Mantena\affilnum{3}, Ruben M. L. Colunga-Biancatelli\affilnum{4,5}}

\affiliation{\affilnum{1}Health, Education, AI and Learning, Office of Enterprise Research and Innovation, Old Dominion University, Norfolk, VA 23508, USA \\
\affilnum{2}National Center for Collaboration in Medical Modeling and Simulation, Macon and Joan Brock Virginia Health Sciences at Old Dominion University, Norfolk, VA 23507, USA \\
\affilnum{3}Eastern Virginia Medical Center, Old Dominion University, Norfolk, VA 23508, USA\\
\affilnum{4}Frank Reidy Research Center for Bioelectrics, Old Dominion University, Norfolk, VA 23508, USA\\
\affilnum{5}Division of Pulmonary and Critical Care Medicine, Department of Medicine, School of Medicine, Macon \& Joan Brock Virginia Health Sciences at Old Dominion University, Norfolk, VA 23507, USA}

\corrauth{Arash Ghorbannia, Old Dominion University, Norfolk, VA, USA.}
\email{[Corresponding author email to be added]}

\begin{abstract}
Nanosecond pulsed electric fields (nsPEF) can stimulate extracellular vesicle (EV) release through coupled membrane electrodynamics, organelle perturbation, Ca$^{2+}$ redistribution, membrane remodeling, repair, vesicle biogenesis, and stress-dependent cargo sorting. However, these processes span nanosecond electrical events, millisecond-to-second Ca$^{2+}$ and organelle signaling, and minute-to-hour EV release and cargo phenotypes, making whole-cell simulation impractical for early model development. Here we define a modular, cell-type-adaptable systems-biology framework for nsPEF-enabled EV engineering. The framework organizes the mechanism into linked layers: pulse delivery and dosimetry; plasma membrane and organelle electrodynamics; ion transport, Ca$^{2+}$ mobilization, ROS, and bioenergetics; Ca$^{2+}$-dependent remodeling and repair; EV biogenesis and subtype release; cargo sorting, EV composition, and potency; injury, debris, and quality gating; and manufacturing, isolation, and quality-control interfaces. Cell-state and disease-state modifiers act as cross-cutting parameter layers that adapt the shared model to cancer cells, endothelial cells, immune cells, stromal cells, and other EV-producing systems. This structure provides a transparent framework for connecting pulse parameters to EV yield, subtype distribution, cargo state, functional potency, purity, and cell viability.
\end{abstract}

\keywords{nanosecond pulsed electric fields; electro-exocytosis; extracellular vesicles; EV engineering; systems biology; compartmental modeling; calcium signaling; electroporation; cargo modulation; biomanufacturing}

\maketitle

\section*{Abbreviations and Nomenclature}
Table~\ref{T_core_abbreviations} summarizes the core biological abbreviations used throughout the manuscript. Variable symbols and parameter notation are listed separately in Appendix Table~\ref{T_A17_glossary}.

\begin{table}[t]
\footnotesize\sf
\centering
\caption{Core abbreviations used throughout the manuscript.\label{T_core_abbreviations}}
\rowcolors{2}{gray!10}{white}
\begin{tabular}{ll}
\toprule
\textbf{Abbr.} & \textbf{Definition} \\
\midrule
nsPEF & nanosecond pulsed electric field \\
EV & extracellular vesicle \\
sEV & small extracellular vesicle \\
m/lEV & medium/large extracellular vesicle \\
AB & apoptotic body \\
MVB & multivesicular body \\
ILV & intraluminal vesicle \\
ER & endoplasmic reticulum \\
ROS & reactive oxygen species \\
ATP & adenosine triphosphate \\
PS & phosphatidylserine \\
ESCRT & endosomal sorting complex required for transport \\
QC & quality control \\
\bottomrule
\end{tabular}
\rowcolors{2}{}{}
\end{table}

\section{Introduction}
Nanosecond pulsed electric fields (nsPEF) are ultrashort, high-intensity electric pulses capable of perturbing plasma membranes and intracellular membrane-bound organelles. Unlike conventional electroporation, which typically uses microsecond-to-millisecond pulses that primarily charge the plasma membrane, nsPEF can access intracellular membranes because the pulse duration is comparable to, or shorter than, relevant membrane-charging times \cite{beebe2003nanosecond,vernier2003calcium,semenov2013primary}. This physical distinction allows nsPEF to trigger rapid Ca$^{2+}$ mobilization, organelle stress responses, phospholipid redistribution, cytoskeletal remodeling, and regulated or stress-associated vesicle release \cite{beebe2012transient,pakhomova2013two,muratori2017activation}. nsPEF have emerged as a bioelectric modality in cancer and other biomedical applications, and recent experimental evidence suggests that they can also stimulate robust extracellular vesicle (EV) release, a process referred to here as \textit{electro-exocytosis} \cite{manoochehri2025extracellular}.

EVs are heterogeneous membrane-bound particles involved in intercellular communication, immune modulation, vascular remodeling, tissue repair, and disease progression. In accordance with MISEV2023 nomenclature principles, this manuscript uses the umbrella term EV when vesicle origin has not been experimentally resolved \cite{welsh2024minimal}. Mechanistic terms such as exosome, microvesicle, and apoptotic body are used when referring to biogenesis-specific model pathways. This distinction is important for modeling because measured EV populations often reflect mixtures of small EVs, medium/large EVs, apoptotic bodies, damaged organellar fragments, protein aggregates, and other co-isolated materials. A useful nsPEF-EV model must therefore predict not only total particle release, but also subtype weighting, cargo state, functional potency, purity, and the boundary between controlled EV release and injury-associated debris.

At the cellular level, nsPEF stimulation initiates a cascade that bridges membrane electrodynamics with vesicle biogenesis. Rapid transmembrane-potential changes and nanopore formation enable ionic fluxes, most notably Ca$^{2+}$ influx and redistribution, with patch-clamp and imaging studies showing minute-scale permeability recovery and Ca$^{2+}$ thresholds that depend on pulse width and cell type \cite{pakhomov2007membrane,pakhomov2007long,nesin2012cell,semenov2013primary}. Intracellular membranes such as the ER, mitochondria, endosomes, and MVBs may also be perturbed, creating routes for store release, mitochondrial stress, ROS generation, trafficking changes, and cargo-sorting modulation \cite{semenov2013recruitment,beebe2012transient,napotnik2012mitochondrial}. Local Ca$^{2+}$ signals then activate effectors including annexins, calpain, scramblases, phospholipases, synaptotagmin-like sensors, and SNARE-associated fusion machinery. These signals drive phosphatidylserine externalization, cytoskeletal reorganization, membrane curvature, MVB docking and fusion, microvesicle budding, and membrane repair \cite{muratori2017activation,bhattacharya2022calcium,williams2023annexin,muratori2021role}.

From a computational systems-biology perspective, nsPEF-induced EV engineering is best formulated as a modular pathway-level perturbation rather than as a complete whole-cell simulation. In this framework, the electrical pulse is represented as a controllable external input that modifies plasma membrane and organelle permeability, producing time-dependent Ca$^{2+}$, ROS, ATP, repair, remodeling, trafficking, and injury states. These upstream variables are then mapped to EV release rates, subtype weighting, cargo composition, functional potency, and quality-control outputs. Such a structure preserves mechanistic interpretability while allowing the same core model to be instantiated in cancer cells, immune cells, endothelial cells, mesenchymal stromal cells, or isolated EVs.

The present manuscript therefore focuses on the modeling problem: how to connect an externally defined nsPEF exposure to membrane and organelle perturbation, Ca$^{2+}$-dependent remodeling, repair, EV biogenesis, cargo modulation, injury-quality gating, and manufacturing-facing EV outputs through a transparent and literature-curatable computational framework.

\begin{figure*}[ht]
    \centering
    \includegraphics[width=0.999\textwidth]{figs/fig01.png}
    \caption{\textbf{Integrated mechanisms linking nanosecond pulsed electric field stimulation to extracellular vesicle (EV) engineering.} Nanosecond pulsed electric fields induce membrane and organelle perturbation, Ca$^{2+}$ redistribution, phosphatidylserine externalization, cytoskeletal remodeling, and repair activation. These processes can promote EV release through endosomal/MVB fusion, plasma membrane budding, and injury-associated vesiculation. In an EV-engineering framework, model outputs include EV yield, size/subtype distribution, cargo composition, functional potency, purity, and cell viability.}
    \abbrevnote{EV, extracellular vesicle; ILV, intraluminal vesicle; MVB, multivesicular body; ESCRT, endosomal sorting complex required for transport; VGCC, voltage-gated calcium channel; PLD2, phospholipase D2; ARF6, ADP-ribosylation factor 6; SNARE, soluble N-ethylmaleimide-sensitive factor attachment protein receptor; LAMP1, lysosome-associated membrane protein 1; HIF-1$\alpha$, hypoxia-inducible factor 1-alpha; IP$_3$R, inositol 1,4,5-trisphosphate receptor; PS, phosphatidylserine.}
    \label{fig:electro_exocytosis_framework}
\end{figure*}

\section{Methods}
We formulated nsPEF-enabled EV engineering as a modular systems-biology framework that links pulse exposure to membrane and organelle perturbation, Ca$^{2+}$ signaling, membrane remodeling, repair, vesicle release, cargo state, injury-quality gating, and manufacturing-facing outputs. Each biological process is represented as a model module with defined inputs, state variables, outputs, and candidate mathematical forms. Paired appendix tables \ref{T_A1_bio_dosimetry}--\ref{T_A16_comp_manufacturing} summarize biological evidence and computational implementation options.

\begin{enumerate}
\item[(i)] Define the minimal set of biological and engineering layers required to connect nsPEF exposure to EV yield, subtype distribution, cargo state, potency, purity, and viability.
\item[(ii)] Assign each layer a computational representation appropriate to its dominant timescale, state variables, measurable outputs, and available modeling literature.
\item[(iii)] Separate biological evidence curation from computational model curation for each module.
\item[(iv)] Distinguish cell-based electro-exocytosis from direct post-isolation EV engineering.
\item[(v)] Represent cell type, disease state, activation state, and population heterogeneity as cross-cutting modifiers rather than as a downstream pathway.
\end{enumerate}

\input{content/in_vitro_size_resolved_calibration_methods}
\input{content/method_global_sensitivity_analysis}

\section{Results}
The resulting framework is a cell-type-adaptable module map rather than a single parameterized model. Table~\ref{T_electro_exo_EV} provides the concise navigation map. Table~\ref{T_modeling_framework} defines the mathematical modeling strategy and module-interface logic. The appendix tables then provide the biological and computational evidence-curation structure needed for future model implementation and parameterization.

\subsection*{Preliminary results}
To demonstrate how the current implementation behaves before formal calibration, we generated a set of targeted example simulations aligned to the first five layers of the framework. These examples are not intended as validation studies. Rather, they show that the selected mathematical forms already produce interpretable regime structure: dosimetry responds to conductivity and pulse-train timing, electrodynamics responds to pulse width and cell size, compartmental ion-transport equations generate coupled Ca$^{2+}$/ROS/ATP trajectories, remodeling rules separate buffering from repair-amplification effects, and the EV-release pool model distinguishes productive secretory windows from injury-dominant vesiculation.

\begin{figure}[t]
\centering
\includegraphics[width=\columnwidth]{figs/temperature_profiles.png}
\caption{\textbf{Preliminary Layer 1--2 example behavior.} The dosimetry comparison shows how the lumped thermal framework separates absorbed electrical dose from retained heating: low-conductivity conditions remain weakly heated, while high-conductivity and high-repetition-rate pulse trains accumulate larger temperature rises even under similar nominal amplitudes.}
\label{fig:prelim_dosimetry}
\end{figure}

\begin{figure}[t]
\centering

\includegraphics[width=\columnwidth]{figs/membrane_voltage_by_pulse_width.png}
\caption{\textbf{Preliminary Layer 1--2 example behavior.} The reduced electrodynamics model shows first-order membrane charging toward the Schwan limit, so longer pulses and larger cells reach higher induced membrane voltages and cross a permeability-relevant threshold more readily than ultrashort or smaller-cell cases.}
\label{fig:prelim_electrodynamics}
\end{figure}


The dosimetry example therefore illustrates the value of a selectable Layer 1 representation rather than a single absorbed-energy scalar. The legacy, adiabatic, and lumped-retention variants produce similar exposure ordering but different thermal interpretations, which is exactly the modeling role intended in Table~\ref{T_A2_comp_dosimetry}. The electrodynamics example then shows the consequence of the Layer 2 reduced charging model: pulse width is not merely a protocol label, but an explicit determinant of how close the membrane gets to its Schwan-limit voltage during each pulse. This matters downstream because the same voltage state drives the permeability and pore-density surrogates that feed the Ca$^{2+}$ and organelle modules.

\begin{figure}[t]
\centering
\includegraphics[width=\columnwidth]{figs/layer3_timeseries_comparison.png}
\caption{\textbf{Preliminary Layer 3--4 example behavior.} The compartmental ion-transport/bioenergetics example compares baseline, higher-dose, calcium-limited, and slow-recovery settings; the chosen ODE framework converts the electroporation state into different Ca$^{2+}$ transients, osmotic stress histories, mitochondrial stress, and ROS accumulation.}
\label{fig:prelim_ion_transport}
\end{figure}


\begin{figure}[t]
\centering
\includegraphics[width=\columnwidth]{figs/remodeling_repair_timeseries.png}
\caption{\textbf{Preliminary Layer 3--4 example behavior.} The remodeling/repair example shows how the reduced Ca$^{2+}$-sensor module separates phosphatidylserine exposure, actin disruption, and resealing competence across buffered, calpain-inhibited, and annexin/lysosome-supported regimes.}
\label{fig:prelim_remodeling_repair}
\end{figure}

The Layer 3 example is especially useful because it shows how the compartmental flux-balance formulation behaves under qualitatively different assumptions. Increasing pulse severity amplifies cytosolic Ca$^{2+}$, ROS, and downstream stress, whereas calcium-limited conditions blunt the Ca$^{2+}$ response even when the same electrodynamic trigger is applied. Slowing ion recovery extends osmotic imbalance, which in turn sustains damage pressure despite a similar initial poration event. These behaviors are consistent with the purpose of Table~\ref{T_A6_comp_ca_ros}: the module is meant to translate a brief electroporation event into longer-lived signaling, organelle, and metabolic trajectories.

The Layer 4 example then shows why the remodeling layer is not redundant with Layer 3. In the present Hill-activation and repair-state framework, buffered microdomains reduce submembrane Ca$^{2+}$ amplification and therefore suppress phosphatidylserine exposure and downstream remodeling. Calpain inhibition preferentially reduces actin disruption without removing the upstream Ca$^{2+}$ transient itself, while stronger annexin/lysosomal support preserves a larger resealing state over time. This is the intended function of Table~\ref{T_A8_comp_remodel_repair}: to resolve whether a given Ca$^{2+}$/stress history is interpreted as recoverable remodeling or escalating injury.

\begin{figure}[t]
\centering
\includegraphics[width=\columnwidth]{figs/ev_biogenesis_timeseries.png}
\caption{\textbf{Preliminary Layer 5 example behavior.} The EV-biogenesis comparison shows that the pool-balance release model distinguishes Rab/ESCRT-limited, ceramide-favored secretory, and injury-shifted regimes. In this framework, MVB pool size and secretory bias respond differently from final apoptotic-body output, allowing the model to separate productive small-EV routing from high-stress vesiculation.}
\abbrevnote{EV, extracellular vesicle; MVB, multivesicular body; ESCRT, endosomal sorting complex required for transport.}
\label{fig:prelim_layer5}
\end{figure}

The Layer 5 example is particularly important because it demonstrates the value of the explicit MVB/ILV/docking/fusion/budding state structure. Rab/ESCRT limitation reduces small-EV release without requiring a generic global suppression of all vesiculation, whereas a ceramide-favored secretory bias increases the release-competent MVB route by reducing lysosomal routing. In contrast, the injury-shift scenario sharply increases apoptotic-body production and collapses viability. This is exactly the interpretive gain intended by the Table~\ref{T_A10_comp_ev_release} framework: subtype release becomes a competition between secretory trafficking, membrane budding, and apoptotic disassembly rather than a single EV-rate output.

\begin{figure*}[t]
\centering
\includegraphics[width=0.98\textwidth]{figs/multilayer_storyboard.png}
\caption{\textbf{Integrated preliminary storyboard across Layers 1--5.} Each row represents one illustrative scenario and each column summarizes one model layer from dosimetry through EV output. The top row shows a mild reversible window, the middle row a more productive secretory window, and the bottom row an injury-dominant window. The figure makes the multiscale logic of the framework visible: stronger Layer 1 dose and Layer 2 charging drive larger Layer 3 Ca$^{2+}$/ROS perturbations, which then reshape Layer 4 repair states and finally alter Layer 5 subtype weighting and viability.}
\abbrevnote{Ca$^{2+}$, calcium ion; ROS, reactive oxygen species; ATP, adenosine triphosphate; EV, extracellular vesicle; sEV, small extracellular vesicle; AB, apoptotic body.}
\label{fig:prelim_storyboard}
\end{figure*}

The integrated storyboard provides a compact narrative view of the framework. In the mild reversible window, the model predicts threshold-level charging, modest temperature rise, limited Ca$^{2+}$ amplification, and predominantly repair-compatible remodeling with low apoptotic output. In the productive secretory window, larger Ca$^{2+}$ loading and stronger repair support coexist with a more favorable secretory routing bias, increasing small-EV output without a catastrophic loss of viability. In the injury-dominant window, the same layered structure yields a different interpretation: stronger charging and thermal load amplify Ca$^{2+}$ and ROS, repair remains incomplete, and apoptotic vesiculation becomes a dominant endpoint. This row-wise transition is encouraging because it shows that the current framework already produces coherent scenario-level behavior while remaining modular enough for later biological calibration.

\begin{figure}[tbp]
\centering
\includegraphics[width=\columnwidth]{figs/downstream_engineering_panel.pdf}
\caption{\textbf{Preliminary Layer 6--8 downstream engineering behavior.} The downstream comparison extends the preliminary examples from vesicle release into product interpretation. A productive cargo window increases protein, RNA, lipid, and antigen cargo scores through secretory and ESCRT-linked terms; a direct-loading mode adds a permeability-driven cargo term while retaining acceptable viability; and an injury-dominant condition may still produce particles but fails the quality gate because apoptosis, necrosis, and contaminant fractions reduce bona fide EV purity and process objective.}
\abbrevnote{EV, extracellular vesicle; RNA, ribonucleic acid; ESCRT, endosomal sorting complex required for transport.}
\label{fig:prelim_downstream_engineering}
\end{figure}

The Layer 6--8 comparison illustrates why the implementation now carries cargo, quality, and manufacturing states beyond total particle release. Productive and direct-loading scenarios can yield similar process objectives through different mechanisms--endogenous cargo sorting versus permeability-assisted loading--whereas the injury-dominant scenario shows that a high particle output is not automatically a useful EV product. This behavior follows the report's computational tables: subtype-weighted cargo and stress sorting feed potency, damage and contaminant mixtures determine the quality gate, and recovery, purity, batch consistency, and potency jointly determine the manufacturing objective.

\input{content/in_vitro_size_resolved_calibration_results}
\input{content/result_global_sensitivity_analysis}

\section{Discussion}
The proposed structure shifts the conceptual endpoint of electro-exocytosis modeling from total EV release alone to nsPEF-enabled EV engineering. In this formulation, EV output is not a single scalar quantity but a multi-dimensional product state that includes particle yield, subtype weighting, cargo composition, potency, purity, and producer-cell viability. This distinction is essential because nsPEF can plausibly increase bona fide EV release, alter cargo sorting, stimulate repair-associated shedding, or generate injury-associated particles depending on pulse dose, cell state, and recovery kinetics. A systems model that does not separate these outputs risks conflating useful EV production with nonspecific stress debris.

A second key feature of the framework is the explicit treatment of organelle electrodynamics. nsPEF effects on the ER, mitochondria, endosomes, and MVBs provide a mechanistic route by which pulse parameters may influence not only Ca$^{2+}$ dynamics but also exosome-biogenesis machinery, cargo sorting, lysosomal routing, and mitochondrial stress. These processes are difficult to resolve experimentally in a single study, but modular modeling allows them to be represented as testable submodels with defined inputs and outputs.

Finally, the framework distinguishes two related but separate applications. In cell-based electro-exocytosis, nsPEF is applied to producer cells to stimulate vesicle release and modify cargo through endogenous signaling and sorting pathways. In direct EV engineering, nsPEF is applied to isolated EVs to transiently permeabilize vesicle membranes for cargo loading or surface modification. These modes share membrane-electroporation principles but require different state variables, size scales, and quality-control outputs. Separating them prevents overextension of a single model while preserving a shared bioelectric modeling backbone.

% ============================================================
% MASTER TABLE: CONCISE NAVIGATION MAP
% ============================================================

\begin{table*}[t]
\footnotesize\sf
\centering
\caption{Concise module map linking nsPEF-enabled EV engineering biology to appendix model-specification tables.\label{T_electro_exo_EV}}
\rowcolors{2}{gray!10}{white}
\begin{tabular}{p{0.7cm} p{3.2cm} p{4.0cm} p{2.6cm} p{2.6cm}}
\toprule
\textbf{Layer} & \textbf{Module} & \textbf{Model interface/output} & \textbf{Biology table} & \textbf{Computational table} \\
\midrule
1 & Pulse delivery, exposure geometry, and dosimetry & $E(t)$, pulse width, pulse number, repetition rate, energy, temperature & Table~\ref{T_A1_bio_dosimetry} & Table~\ref{T_A2_comp_dosimetry} \\
2 & Plasma membrane and organelle electrodynamics & $\Delta V_m$, $\Delta V_{ER}$, $\Delta V_{mito}$, $\Delta V_{MVB}$, $P_i(t)$ & Table~\ref{T_A3_bio_electrodynamics} & Table~\ref{T_A4_comp_electrodynamics} \\
3 & Ion transport, Ca$^{2+}$ mobilization, ROS, and bioenergetics & $J_{\mathrm{Ca}}$, $C_i$, $C_{ER}$, $C_{mito}$, ROS, ATP & Table~\ref{T_A5_bio_ca_ros} & Table~\ref{T_A6_comp_ca_ros} \\
4 & Ca$^{2+}$-dependent remodeling and repair & PS exposure, calpain, annexins, actin state, resealing & Table~\ref{T_A7_bio_remodel_repair} & Table~\ref{T_A8_comp_remodel_repair} \\
5 & EV biogenesis and subtype release & $R_{sEV}$, $R_{m/lEV}$, $R_{AB}$, subtype weighting & Table~\ref{T_A9_bio_ev_release} & Table~\ref{T_A10_comp_ev_release} \\
6 & Cargo sorting, EV composition, and potency & protein/RNA/lipid cargo, antigen enrichment, functional potency & Table~\ref{T_A11_bio_cargo} & Table~\ref{T_A12_comp_cargo} \\
7 & Injury, debris, and quality gate & viability, apoptosis, necrosis, debris fraction, contamination markers & Table~\ref{T_A13_bio_quality_gate} & Table~\ref{T_A14_comp_quality_gate} \\
8 & Manufacturing, isolation, and QC interface & yield, recovery, purity, batch consistency, scalability & Table~\ref{T_A15_bio_manufacturing} & Table~\ref{T_A16_comp_manufacturing} \\
Cross-cutting & Cell-state and disease modifiers & parameter scalars, phenotype priors, heterogeneity & Tables~\ref{T_A1_bio_dosimetry}--\ref{T_A15_bio_manufacturing} & Tables~\ref{T_A2_comp_dosimetry}--\ref{T_A16_comp_manufacturing} \\
\bottomrule
\end{tabular}
\vspace{0.5em}
\rowcolors{2}{}{}
\abbrevnote{nsPEF, nanosecond pulsed electric field; EV, extracellular vesicle; sEV, small extracellular vesicle; m/lEV, medium/large extracellular vesicle; AB, apoptotic body; PS, phosphatidylserine; ROS, reactive oxygen species; QC, quality control.}
\end{table*}

% ============================================================
% MAIN TEXT MODELING FRAMEWORK TABLE
% ============================================================

\begin{table*}[t]
\footnotesize\sf
\centering
\caption{Modeling framework and module-interface logic for the integrated nsPEF-EV engineering model.\label{T_modeling_framework}}
\rowcolors{2}{gray!10}{white}
\begin{tabular}{p{2.8cm} p{3.3cm} p{3.1cm} p{3.1cm} p{3.1cm}}
\toprule
\textbf{Module} & \textbf{Mathematical representation} & \textbf{Inputs} & \textbf{Outputs} & \textbf{Rationale} \\
\midrule
Pulse delivery and dosimetry & External forcing, exposure geometry, energy/thermal constraints & pulse protocol, chamber geometry, conductivity, cell density & $E(t)$, dose metrics, temperature & Defines controllable engineering input and batch-to-batch exposure. \\
Membrane and organelle electrodynamics & Circuit model plus pore-density/pore-radius kinetics & $E(t)$, geometry, membrane properties, organelle size & $\Delta V_m$, $\Delta V_o$, $G_m$, $P_i$ & Converts nsPEF exposure into plasma-membrane and organelle permeability. \\
Ion transport, Ca$^{2+}$, ROS, ATP & Compartmental flux-balance ODEs with optional GHK/Nernst--Planck fluxes & permeability, ion gradients, pumps, channels, stores & $C_i$, $C_{sub}$, $C_{ER}$, $C_{mito}$, ROS, ATP & Links electroporation to signaling, organelle stress, and recovery capacity. \\
Remodeling and repair & Reaction--diffusion, buffer kinetics, Hill activation, repair ODEs & Ca$^{2+}$ transient, ROS, membrane damage & effector activation, PS exposure, actin state, resealing & Defines reversible remodeling, budding competence, and membrane recovery. \\
EV release & Pool-balance and Ca$^{2+}$-dependent rate models & MVB pool, budding sites, Ca$^{2+}$, PS, cytoskeletal state & $R_{sEV}$, $R_{m/lEV}$, $R_{AB}$ & Maps molecular states to size/subtype-resolved vesicle release. \\
Cargo and potency & Cargo-sorting weights, reaction networks, or phenotype-weighted output models & ESCRT, Rab, tetraspanins, stress state, EV subtype & protein/RNA/lipid cargo, potency score & Captures EV engineering beyond yield alone. \\
Injury and quality gate & Threshold rules, hybrid state transitions, contamination models & residual conductance, Ca$^{2+}$, ATP, ROS, caspases & viable/stressed/apoptotic fractions, debris burden & Separates useful EV release from injury-associated particles. \\
Manufacturing and QC & Process-yield balances, recovery models, mixed-effects batch models & harvest time, isolation method, QC markers, scale & yield, recovery, purity, reproducibility & Connects cell-level biology to translational EV production. \\
Cell-state modifiers & Parameter scaling, priors, mixed-effects or ensemble models & cell type, disease state, activation, stiffness, lipid state & module-specific sensitivity factors & Adapts the shared model to different producer cells and disease contexts. \\
\bottomrule
\end{tabular}
\vspace{0.5em}
\rowcolors{2}{}{}
\footnotesize
\textbf{Note:} The table defines modeling form at the module level. Appendix tables provide submodule-specific biological evidence and computational implementation options.
\abbrevnote{nsPEF, nanosecond pulsed electric field; EV, extracellular vesicle; MVB, multivesicular body; ESCRT, endosomal sorting complex required for transport; PS, phosphatidylserine; ROS, reactive oxygen species; ATP, adenosine triphosphate; QC, quality control; GHK, Goldman-Hodgkin-Katz.}
\end{table*}

%\begin{acks}
% Funding and acknowledgements to be added.
%\end{acks}

\section*{Appendix}
\appendix

\renewcommand{\theequation}{A\arabic{equation}}
\renewcommand{\thefigure}{A\arabic{figure}}
\renewcommand{\thetable}{A\arabic{table}}

\setcounter{equation}{0}
\setcounter{figure}{0}
\setcounter{table}{0}

\subsection*{Appendix organization and model-curation strategy}
The appendix is organized as paired biological and computational tables for each modeling layer. The biological tables document the experimental and mechanistic rationale for including each submodule. The computational tables document candidate mathematical representations, variables, inputs, outputs, and parameterization strategies. Evidence confidence is determined by literature review using the following working categories: established, supported, hypothesized, or model assumption. Table pairs are referenced locally in each module subsection so that biological claims and candidate model forms can be traced back to their evidence and implementation tables.

% ============================================================
\subsection*{Pulse Delivery, Exposure Geometry, and Dosimetry}
% ============================================================

Pulse delivery, exposure geometry, and dosimetry define the engineered exposure that initiates electro-exocytosis. In this layer, the pulse waveform denotes the externally imposed electric-field history, including pulse amplitude, duration, number, polarity, and repetition rate. Early nsPEF studies show that these variables set the transition from plasma-membrane poration to intracellular Ca$^{2+}$ mobilization, PS exposure, and cell-fate responses \cite{beebe2003nanosecond,vernier2003calcium,pakhomov2007membrane,semenov2013primary}. Exposure geometry denotes the physical arrangement of electrodes, chamber, cuvette, culture format, or flow path that converts a programmed voltage waveform into the electric field experienced by the sample. Because cells may receive different effective exposures depending on electrode gap, edge effects, medium conductivity, cell density, sample volume, and position within the device, geometry is central to reproducibility \cite{joshi2000electroporation,kotnik2019membrane,baker2024numerical}. A nominal pulse protocol is therefore not sufficient unless the exposure configuration, field uniformity, and measurement method are also defined \cite{rems2013cell,leveque2017measurement}. The biological/process evidence for these exposure variables is summarized in Table~\ref{T_A1_bio_dosimetry}, and the corresponding computational abstractions are listed in Table~\ref{T_A2_comp_dosimetry}.

For the present compartmental modeling framework, the spatial field calculation should be treated as an upstream dosimetry and calibration step rather than as a spatially resolved biological state. Device settings and exposure geometry can be collapsed into lumped descriptors such as an effective field amplitude, a cumulative pulse-dose index, a heating constraint or temperature history, an exposure-uniformity factor, and the recovery or harvest interval. These descriptors define controllable engineering inputs and batch-to-batch exposure conditions for the downstream ODE modules, while avoiding the need to track chamber position as a dynamic state variable. This separation is also biologically important: energy deposition and Joule heating must be monitored so that field-driven signaling can be distinguished from thermal artifacts, and the recovery interval must be specified because EV yield, cargo composition, and viability are measured over minutes to hours after pulsing rather than at the instant of exposure. Recent nsPEF spheroid work explicitly links pulse width, repetition rate, buffer conductivity, absorbed energy density, and measured temperature rise, showing that high-conductivity media and compressed pulse trains can produce substantially different thermal histories even when nominal pulse amplitudes are similar \cite{orlacchio2023effects}. Related cancer-ablation studies emphasize the same practical distinction between absorbed electrical dose and heat retained during the treatment interval \cite{gudvangen2022electroporation}. Figure~\ref{fig:dosimetry_schematic} summarizes this interface between device-level pulse delivery and reduced model inputs.

We therefore use a tiered dosimetry strategy in which increasingly detailed pulse-energy models can be selected according to the available calibration data. The simplest baseline preserves a legacy absorbed-energy descriptor,
\begin{equation}
u_{\mathrm{legacy}}=\sigma E_{\mathrm{peak}}^{2}t_pN ,
\end{equation}
where $\sigma$ is the medium conductivity, $E_{\mathrm{peak}}$ is the nominal peak electric field, $t_p$ is pulse width, and $N$ is pulse number. This expression is useful as a transparent dose index and maintains continuity with early exploratory simulations, but it does not distinguish square, bipolar, or exponentially decaying waveforms and treats geometry correction only as a coarse scalar.

A waveform-aware Joule model retains the same physical skeleton but weights the pulse by the normalized squared waveform integral,
\begin{align}
\chi_w &= \frac{1}{t_p}\int_0^{t_p}\left(\frac{E(t)}{E_{\mathrm{peak}}}\right)^2dt,\\
u_J &= \sigma E_{\mathrm{peak}}^{2}t_pN\chi_w\langle g_E^2\rangle,\\
\Delta T_{\mathrm{ad}} &= \frac{u_J}{C_v}.
\end{align}
Here $\chi_w$ captures waveform shape, $\langle g_E^2\rangle$ represents the geometry-averaged squared-field factor, and $C_v$ is the volumetric heat capacity. This adiabatic estimate is intentionally conservative: it represents the temperature rise that would occur if deposited energy remained locally retained. It is appropriate for ranking protocols by electrical energy deposition, but it should not be interpreted as a measured sample temperature when heat diffusion, convection, electrode heat exchange, or long inter-pulse intervals are relevant.

When temperature history is needed, the adiabatic estimate can be coupled to a reduced thermal-retention model,
\begin{equation}
\Delta T_{\mathrm{end}}=\eta\Delta T_{\mathrm{ad}}
\frac{\tau_{\mathrm{loss}}}{T_{\mathrm{train}}}
\left(1-e^{-T_{\mathrm{train}}/\tau_{\mathrm{loss}}}\right),
\end{equation}
where $T_{\mathrm{train}}\approx N/f_{\mathrm{rep}}$ is the pulse-train duration, $\tau_{\mathrm{loss}}$ is an empirical thermal relaxation time, and $\eta$ is a setup-dependent retention factor. This reduced form is not a substitute for finite-element thermal modeling, but it captures the main experimental observation that identical absorbed energy can lead to different end-of-train temperature rises when pulse repetition rate, train duration, buffer conductivity, and exposure geometry differ \cite{orlacchio2023effects,gudvangen2022electroporation}. Because temperature can itself modulate membrane fluidity, permeabilization thresholds, and cell fate, thermal terms should be reported as independent dosimetry outputs rather than folded silently into an electrical dose scalar \cite{kanduser2008temperature,song2011synergistic,song2017thermal,yan2021synergistic,pakhomova2013two}.

\begin{figure*}[t]
    \centering
    \includegraphics[width=0.95\textwidth]{figs/figA1_dosimetry_schematic.png}
    \caption{\textbf{Conceptual interface for pulse delivery, exposure geometry, and dosimetry.} The programmed nsPEF waveform and physical exposure configuration determine the actual sample exposure. In a reduced compartmental framework, spatial field information is used for upstream calibration and then summarized as lumped descriptors such as effective field amplitude, cumulative pulse dose, heating constraint, exposure-uniformity factor, and recovery interval. These descriptors provide reproducible inputs to downstream biological ODE modules without requiring a spatially resolved core model.}
    \abbrevnote{nsPEF, nanosecond pulsed electric field; ODE, ordinary differential equation.}
    \label{fig:dosimetry_schematic}
\end{figure*}

\begin{table*}[t]
\footnotesize\sf
\centering
\caption{Biological evidence table for pulse delivery, exposure geometry, and dosimetry.\label{T_A1_bio_dosimetry}}
\rowcolors{2}{gray!10}{white}
\begin{tabular}{p{2.4cm} p{3.0cm} p{3.8cm} p{2.7cm} p{3.4cm}}
\toprule
\textbf{Submodule} & \textbf{Biological/process phenomenon} & \textbf{Evidence to curate} & \textbf{Relevance to EV engineering} & \textbf{Confidence} \\
\midrule
Pulse waveform & nsPEF duration, amplitude, number, repetition rate & Parameter-response studies across pulse width, field strength, repetition, and pulse number & Defines controllable input space for EV yield and cargo modulation & Established \cite{beebe2003nanosecond,vernier2003calcium,pakhomov2007membrane,semenov2013primary,orlacchio2023effects} \\
Exposure geometry & Electrode/chamber/cuvette/flow geometry & Field uniformity, edge effects, cell density, medium conductivity & Controls reproducibility and batch exposure & Supported; implementation remains lumped \cite{joshi2000electroporation,kotnik2019membrane,baker2024numerical,leveque2017measurement,orlacchio2023effects} \\
Energy and heating & Joule heating and thermal artifact risk & Absorbed energy density, waveform energy factor, temperature traces, duty-cycle and repetition-rate effects & Separates bioelectric effects from heat-induced stress & Supported; setup-specific calibration needed \cite{orlacchio2023effects,gudvangen2022electroporation,kanduser2008temperature,song2011synergistic,song2017thermal} \\
Recovery interval & Post-pulse harvest window & EV release and repair kinetics over seconds to hours & Determines when yield, cargo, and viability are measured & Supported\cite{silkunas2026peculiarities,szwedowicz2026reversible,bhattacharya2022calcium,muratori2017activation,williams2023annexin} \\
Scale-up mode & Static versus flow-through treatment & Flow exposure, residence time, process monitoring, and sterility constraints & Connects mechanism to manufacturing implementation & Hypothesized\cite{liu2026microfluidic,bakute2024microphysiological,garcia2016microfluidic,bonakdar2017microfluidic,li2015electroporation} \\
\bottomrule
\end{tabular}
\rowcolors{2}{}{}
\abbrevnote{nsPEF, nanosecond pulsed electric field; EV, extracellular vesicle.}
\end{table*}

\begin{table*}[t]
\footnotesize\sf
\centering
\caption{Computational framework table for pulse delivery, exposure geometry, and dosimetry.\label{T_A2_comp_dosimetry}}
\rowcolors{2}{gray!10}{white}
\begin{tabular}{p{2.6cm} p{3.0cm} p{4.6cm} p{2.8cm} p{2.8cm}}
\toprule
\textbf{Submodule} & \textbf{Model form} & \textbf{Representative relation} & \textbf{Inputs} & \textbf{Outputs} \\
\midrule
Pulse input & External forcing function \cite{beebe2003nanosecond,vernier2003calcium,pakhomov2007membrane,semenov2013primary} & $E(t)=E_0u(t;t_p,N,f)$ & amplitude, width, number, repetition & $E(t)$ \\
Dose metric & Selectable absorbed-energy model \cite{orlacchio2023effects,gudvangen2022electroporation} & $u=\sigma E_{\mathrm{peak}}^2t_pN\chi_w\langle g_E^2\rangle$ & pulse train, conductivity, waveform, geometry & energy density, dose index \\
Field distribution & Laplace/quasi-static field model \cite{joshi2000electroporation,kotnik2019membrane,baker2024numerical,leveque2017measurement} & $\nabla\cdot(\sigma\nabla\phi)=0$, $E=-\nabla\phi$ & chamber geometry, conductivity & local $E(\mathbf{x},t)$ \\
Thermal constraint & Adiabatic or lumped heat-loss model \cite{orlacchio2023effects,gudvangen2022electroporation,kanduser2008temperature,song2011synergistic,song2017thermal} & $\Delta T_{\mathrm{end}}=F_{\mathrm{ret}}\Delta T_{\mathrm{ad}}$ & field, medium, cooling, pulse train & $T(t)$, $\Delta T_{\mathrm{end}}$ \\
Flow-through exposure & Residence-time model \cite{liu2026microfluidic,garcia2016microfluidic,bonakdar2017microfluidic,li2015electroporation} & $D_E(\mathbf{x})=\int_{\tau} E(\mathbf{x}(t))^2dt$ & flow rate, chamber field & exposure distribution \\
\bottomrule
\end{tabular}
\rowcolors{2}{}{}
\abbrevnote{EV, extracellular vesicle.}
\end{table*}
% ============================================================
\subsection*{Plasma Membrane and Organelle Electrodynamics}
% ============================================================

A central goal of the framework is to connect nsPEF exposure to downstream Ca$^{2+}$-dependent EV release and cargo modulation without requiring a complete whole-cell reconstruction. The membrane and organelle-electrodynamics layer therefore serves as the upstream bioelectric perturbation module. The applied pulse induces transient plasma-membrane polarization, electropore formation, increased conductance, and ion-specific permeability, all of which have been observed or modeled in single-cell electroporation studies \cite{kotnik2000analytical,krassowska2007modeling,pakhomov2007long}. Short pulses can also couple to intracellular membranes, supporting reduced organelle-voltage terms for ER, mitochondrial, endosomal, and MVB-related perturbation \cite{vernier2003calcium,semenov2013primary,beebe2012transient,baker2024numerical}. These outputs provide the primary interface between nsPEF exposure and downstream Ca$^{2+}$, repair, exocytosis, EV-biogenesis, and cargo-sorting modules.

Computational electroporation literature clarifies how much of this upstream module can remain reduced while still being mechanistically anchored. Molecular dynamics studies show field-driven water defects progressing into hydrophilic pores on nanosecond time scales, with pore formation, water-channel growth, voltage removal, and resealing all occurring as distinct processes \cite{tieleman2004molecular,tarek2005membrane}. Vesicle electromechanics adds the continuum constraint that breakdown voltage depends on membrane tension, area expansivity, and lipid composition \cite{needham1989electro}. Together, these papers support the present choice to keep pore nucleation, pore-radius/tension effects, permeability, and resealing as separable calibrated states rather than one binary permeabilization switch. Table~\ref{T_A3_bio_electrodynamics} summarizes the biological electrodynamics evidence, and Table~\ref{T_A4_comp_electrodynamics} lists the reduced computational forms used for this layer.

The current implementation uses a reduced electrodynamics ladder rather than a full spatial pore-kinetic simulation. Plasma-membrane polarization is represented by a Schwan-style limiting voltage scaled by the effective field, cell radius, and a cell-state membrane modifier, with a first-order charging correction for finite pulse width:
\[
\begin{aligned}
\Delta V_m^{\infty} &= f_sE_{\mathrm{eff}}aM_m, \\
\Delta V_m(t_p) &= \Delta V_m^{\infty}
\left(1-\exp[-t_p/\tau_m]\right).
\end{aligned}
\]
where $f_s$ is the Schwan geometry factor, $E_{\mathrm{eff}}$ is the geometry-corrected electric field, $a$ is the configurable cell radius, $M_m$ is the membrane-state modifier, $t_p$ is pulse width, and $\tau_m$ is an effective membrane charging time \cite{kotnik2000analytical,joshi2000electroporation,kotnik2019membrane}. This formulation resolves the previous fixed-radius placeholder by making $\Delta V_m$ explicitly cell-size dependent; larger cells reach higher induced membrane voltages at the same field, while very short pulses may terminate before the membrane reaches its Schwan limit.

Electroporation is then represented by reduced permeability and pore-density closures:
\[
\begin{aligned}
P_m &= \left[1+\exp[-s(\Delta V_m-V_{\mathrm{th}})]\right]^{-1}, \\
N_p &= N_{\max}
\frac{\Delta V_m^n}{V_{1/2}^n+\Delta V_m^n}.
\end{aligned}
\]
These terms are phenomenological surrogates for voltage-dependent pore creation, pore growth, conductance, and resealing models \cite{vasilkoski2006absolute,krassowska2007modeling}. Patch-clamp studies provide direct calibration targets because nsPEF-induced conductance can persist for minutes and recover in a cell-type-dependent manner \cite{pakhomov2007membrane,pakhomov2007long,nesin2012cell}. The reduced pore-density variable should therefore be interpreted as an effective permeability state, not an absolute physical pore count, until it is calibrated against conductance, dye uptake, ion flux, or recovery data. Organelle perturbation is currently represented by configurable fractions of the plasma-membrane voltage, $\Delta V_o=\alpha_o\Delta V_m$ for ER, mitochondria, and MVB/endosomal membranes. This preserves an explicit route from nsPEF exposure to intracellular membrane stress while marking the organelle layer as calibration-dependent until multi-shell, geometry-resolved, or organelle-specific simulations are incorporated \cite{semenov2013primary,beebe2012transient,baker2024numerical}. The same pore state also feeds the lipid-remodeling layer because experiments and MD simulations connect nanosecond-pulse nanopores with voltage-driven phosphatidylserine translocation \cite{vernier2006nanoporeJACS,vernier2006nanoporePhysBio}.

\begin{table*}[t]
\footnotesize\sf
\centering
\caption{Biological evidence table for plasma membrane and organelle electrodynamics.\label{T_A3_bio_electrodynamics}}
\rowcolors{2}{gray!10}{white}
\begin{tabular}{p{2.4cm} p{3.1cm} p{3.7cm} p{2.8cm} p{4.0cm}}
\toprule
\textbf{Submodule} & \textbf{Biological phenomenon} & \textbf{Evidence to curate} & \textbf{Relevance to EV engineering} & \textbf{Confidence} \\
\midrule
Membrane polarization & Rapid transmembrane-potential rise & Measurements or simulations of $\Delta V_m$ under nsPEF exposure & Determines whether local membrane regions reach electroporation range & Established; reduced Schwan/charging implementation \cite{kotnik2000analytical,joshi2000electroporation,kotnik2019membrane,baker2024numerical} \\
Plasma-membrane nanopores & Transient nanoscale membrane defects & Conductance, dye uptake, ion permeability, recovery kinetics & Creates pathways for Ca$^{2+}$ and other ion entry & Established; pore-density proxy remains phenomenological \cite{pakhomov2007membrane,pakhomov2007long,vasilkoski2006absolute,krassowska2007modeling,kotnik2019membrane} \\
Organelle electroporation & ER, mitochondrial, nuclear, endosomal, and MVB membrane perturbation & Ca$^{2+}$ release, mitochondrial depolarization, intracellular membrane access & Provides route for store release, ROS, and altered MVB/cargo dynamics & Supported; organelle fractions need calibration \cite{vernier2003calcium,semenov2013primary,beebe2012transient,baker2024numerical} \\
Permeability recovery & Conductance decay and resealing & Patch clamp, uptake assays, dye exclusion & Sets duration and magnitude of post-pulse ion flux & Established; kinetics are cell-type dependent \cite{pakhomov2007long,nesin2012cell,pakhomov2007membrane,pakhomova2013two} \\
Isolated EV permeabilization & Direct nsPEF effect on EV membranes & Vesicle integrity, cargo loading, aggregation, marker retention & Defines second mode for post-isolation EV engineering & Hypothesized; not implemented in the current cell-based layer \cite{singh2025electroporation,hood2014maximizing,kooijmans2013electroporation,fuhrmann2015active,didiot2018optimized} \\
\bottomrule
\end{tabular}
\rowcolors{2}{}{}
\abbrevnote{nsPEF, nanosecond pulsed electric field; ER, endoplasmic reticulum; EV, extracellular vesicle; MVB, multivesicular body.}
\end{table*}

\begin{table*}[t]
\footnotesize\sf
\centering
\caption{Computational framework table for plasma membrane and organelle electrodynamics.\label{T_A4_comp_electrodynamics}}
\rowcolors{2}{gray!10}{white}
\begin{tabular}{p{2.6cm} p{3.0cm} p{4.6cm} p{2.8cm} p{2.8cm}}
\toprule
\textbf{Submodule} & \textbf{Model form} & \textbf{Representative relation} & \textbf{Inputs} & \textbf{Outputs} \\
\midrule
Membrane charging & Reduced Schwan/circuit model \cite{kotnik2000analytical,joshi2000electroporation,kotnik2019membrane} & $\Delta V_m=f_sE_{\mathrm{eff}}a(1-e^{-t_p/\tau_m})$ & $E(t)$, geometry, cell radius, membrane state & $\Delta V_m$ \\
Organelle charging & Reduced multi-shell surrogate \cite{semenov2013primary,beebe2012transient,baker2024numerical} & $\Delta V_o=\alpha_o\Delta V_m$ & organelle factors, plasma-membrane voltage & $\Delta V_o$ \\
Pore nucleation & Voltage-dependent permeability proxy \cite{tieleman2004molecular,tarek2005membrane} & $P_m=[1+\exp[-s(\Delta V_m-V_{\mathrm{th}})]]^{-1}$ & $\Delta V_m$, threshold parameters & $P_m$ \\
Pore-radius evolution & Pore-energy or reduced kinetic model \cite{needham1989electro,tarek2005membrane} & $\frac{dr_p}{dt}=F(r_p,\Delta V,\gamma,\sigma)$ & pore energy, tension & $r_p(t)$ \\
Permeability closure & Saturating pore-density closure \cite{smith2004stable,krassowska2007modeling} & $N_p=N_{\max}\Delta V_m^n/(V_{1/2}^n+\Delta V_m^n)$ & pore-density scale, voltage scale, Hill exponent & $N_p$, $P_i$ \\
Isolated EV charging & Vesicle-scale circuit model \cite{singh2025electroporation,hood2014maximizing,kooijmans2013electroporation} & $C_{EV}\frac{d\Delta V_{EV}}{dt}+G_{EV}\Delta V_{EV}=f(E,r_{EV})$ & EV radius, medium conductivity & $P_{EV}(t)$, integrity state \\
\bottomrule
\end{tabular}
\rowcolors{2}{}{}
\abbrevnote{ER, endoplasmic reticulum; EV, extracellular vesicle; MVB, multivesicular body.}
\end{table*}

% ============================================================
\subsection*{Ion Transport, Ca$^{2+}$ Mobilization, ROS, and Bioenergetics}
% ============================================================

The ion-transport and organelle-state module translates nsPEF-induced membrane permeability into intracellular ionic, metabolic, and stress responses. The minimal model treats electropore-mediated Ca$^{2+}$ influx and ER release as source terms in a cytosolic Ca$^{2+}$ balance. More detailed formulations include voltage-gated or leak-channel entry, ER release through IP$_3$ receptors or ryanodine receptors, SERCA-mediated refilling, PMCA and NCX extrusion, mitochondrial uptake and release, ROS generation, ATP dependence, and osmotic recovery. Table~\ref{T_A5_bio_ca_ros} organizes the biological evidence for these submodules, while Table~\ref{T_A6_comp_ca_ros} lists the corresponding reduced flux-balance and organelle-state model forms.

For the current computational framework, this layer is represented as a reduced flux-balance ODE system rather than a full electrodiffusion model. The pore term is treated as a transient extracellular Ca$^{2+}$ source gated by the electrodynamic permeability state. This choice is supported by studies in which nsPEF rapidly elevates cytosolic Ca$^{2+}$ and the response depends on extracellular Ca$^{2+}$, pulse width, and membrane pathway \cite{vernier2003calcium,bagalkot2018membrane,yun2024modulating}. In CHO cells, a 60-ns pulse produced separable external-Ca$^{2+}$ and ER-release thresholds and a critical Ca$^{2+}$ range for amplification, making these data especially useful for tuning reduced Ca$^{2+}$ source terms \cite{semenov2013primary}.

ER mobilization is modeled as an electrically triggered store-release pathway with SERCA refilling, while IP$_3$R/RyR-mediated amplification is retained as a cell-type-specific extension for excitable or store-sensitive cells \cite{semenov2013recruitment,tolstykh2013activation,tolstykh2014pip2,wang2009cardiac,azarov2019excitation}. The computational calcium literature gives useful bounds for that extension: Li and Rinzel reduced a detailed IP$_3$ receptor model to a Hodgkin--Huxley-like form with IP$_3$ regulation, Ca$^{2+}$ activation, and Ca$^{2+}$ inactivation, while Shannon and coworkers used junctional, subsarcolemmal, and bulk compartments with realistic buffering, reversible SR uptake, Na$^+$--Ca$^{2+}$ exchange, and load-dependent release \cite{li1994equations,shannon2004mathematical}. We therefore keep the default ER term compact, but expose channel opening, store content, pump capacity, and exchanger recovery as parameters when cell-type-specific data are available.

Na$^+$, K$^+$, and Cl$^-$ fluxes are included as lumped electropore/channel-mediated perturbations that relax toward baseline and drive an osmotic-stress state. This abstraction is justified by evidence that nsPEF increases leak conductance, depolarizes membrane potential, alters voltage-gated channels, and changes cell volume \cite{pakhomov2007membrane,smith2011transmembrane,burke2017nanosecond,yang2021regulation}. Na$^+$/K$^+$-ATPase sensitivity is retained as a recovery modifier because ATP1A1 silencing attenuates membrane disruption after nanosecond pulses \cite{silkuniene2023silencing}.

Mitochondrial Ca$^{2+}$ uptake/release and depolarization connect cytosolic Ca$^{2+}$ to ROS and ATP availability. nsPEF studies support mitochondrial membrane perturbation and ROS generation, while broader calcium-electroporation data show that Ca$^{2+}$ overload can deplete ATP and shift the cell toward injury \cite{beebe2012transient,napotnik2012mitochondrial,pakhomova2012oxidative,nuccitelli2013reactive,radzeviciute2024calcium,gibot2020calcium}. On the computational side, mitochondrial modules can be scaled from a minimal Ca$^{2+}$-handling block with respiration, ATP synthase, proton leak, adenine nucleotide exchange, uniporter, and Na$^+$--Ca$^{2+}$ exchange \cite{magnus1997minimal}, to a simplified ATP-production model driven by glycolytic and Ca$^{2+}$ inputs \cite{bertram2006simplified}, or to a thermodynamically balanced oxidative-phosphorylation model that conserves mass, charge, and free energy \cite{beard2005biophysical}. ROS is represented with Ca$^{2+}$-, mitochondrial-, osmotic-, and pulse-dependent source terms plus first-order clearance; this reduced form is consistent with models in which ROS production, scavenging, and ROS-sensitive inner-membrane conductance shape mitochondrial potential, NADH, and redox dynamics \cite{cortassa2004mitochondrial,murphy2016mitochondrial,zhang2022crosstalk}. ATP is kept as a normalized production-consumption state coupled to mitochondrial potential, pump activity, repair demand, and oxidative stress.

\begin{equation}
\begin{aligned}
\frac{d[\mathrm{Ca}^{2+}]_i}{dt} &=
J_{\mathrm{Ca,pore}} + J_{\mathrm{VGCC}} + J_{\mathrm{ER,rel}} - J_{\mathrm{SERCA}} \\
&\quad - J_{\mathrm{PMCA}} - J_{\mathrm{NCX}} - J_{\mathrm{mito}} - J_{\mathrm{buffer}} .
\end{aligned}
\end{equation}

\begin{table*}[t]
\footnotesize\sf
\centering
\caption{Biological evidence table for ion transport, Ca$^{2+}$ mobilization, ROS, and bioenergetics.\label{T_A5_bio_ca_ros}}
\rowcolors{2}{gray!10}{white}
\begin{tabular}{p{2.7cm} p{3.4cm} p{4.2cm} p{3.0cm} p{2.0cm}}
\toprule
\textbf{Submodule} & \textbf{Biological phenomenon} & \textbf{Evidence to curate} & \textbf{Relevance to EV engineering} & \textbf{Confidence} \\
\midrule
Electropore Ca$^{2+}$ influx & Rapid extracellular Ca$^{2+}$ entry & Ca$^{2+}$ imaging, chelation, Ca$^{2+}$-free media & Primary trigger for downstream release and repair mechanisms & Established \cite{bagalkot2018membrane,bhattacharya2022calcium,morotomi2014extracellular,yun2024modulating,cantu2021caveolin} \\
ER Ca$^{2+}$ release & Intracellular store mobilization & IP$_3$R/RyR involvement, ER store depletion & Amplifies cytosolic Ca$^{2+}$ transients and fusion signals & Established \cite{semenov2013recruitment,tolstykh2013activation,tolstykh2014pip2,wang2009cardiac,azarov2019excitation} \\
Na$^+$, K$^+$, Cl$^-$ fluxes & Ionic/osmotic perturbation & Membrane-potential recovery, swelling, ionic imbalance & Modulates volume, tension, and repair state & Supported \cite{smith2011transmembrane,burke2017nanosecond,muratori2017activation,silkuniene2023silencing,yang2021regulation} \\
Mitochondrial Ca$^{2+}$ & Uptake/release and depolarization & Mitochondrial Ca$^{2+}$, $\Delta\Psi_m$, mPTP evidence & Links Ca$^{2+}$ to metabolic stress and fate decisions & Supported \cite{beebe2012transient,napotnik2012mitochondrial,murphy2016mitochondrial} \\
ROS generation & Oxidative stress after nsPEF & ROS probes, antioxidant rescue, mitochondrial source evidence & Modifies EV secretion, cargo, and injury transition & Supported \cite{pakhomova2012oxidative,nuccitelli2013reactive,zhang2022crosstalk,lee2019caveolin,wiczew2021molecular} \\
ATP/bioenergetics & Energy availability for pumps and trafficking & ATP depletion, pump recovery, metabolic assays & Controls repair, cargo sorting, and EV-release capacity & Hypothesized \cite{frandsen2012calcium,radzeviciute2024calcium,gibot2020calcium} \\
\bottomrule
\end{tabular}
\rowcolors{2}{}{}
\abbrevnote{EV, extracellular vesicle; ER, endoplasmic reticulum; IP$_3$R, inositol 1,4,5-trisphosphate receptor; RyR, ryanodine receptor; ROS, reactive oxygen species; ATP, adenosine triphosphate; mPTP, mitochondrial permeability transition pore.}
\end{table*}

\begin{table*}[t]
\footnotesize\sf
\centering
\caption{Computational framework table for ion transport, Ca$^{2+}$ mobilization, ROS, and bioenergetics.\label{T_A6_comp_ca_ros}}
\rowcolors{2}{gray!10}{white}
\begin{tabular}{p{2.6cm} p{3.0cm} p{4.6cm} p{2.8cm} p{2.8cm}}
\toprule
\textbf{Submodule} & \textbf{Model form} & \textbf{Representative relation} & \textbf{Inputs} & \textbf{Outputs} \\
\midrule
Pore-mediated Ca$^{2+}$ flux & GHK/Ohmic flux \cite{tieleman2004molecular,tarek2005membrane} & $J_{\mathrm{Ca,pore}}=P_{\mathrm{Ca}}(t)\Phi_{\mathrm{Ca}}(C_e,C_i,V_m)$ & $P_{\mathrm{Ca}}$, gradients, voltage & Ca$^{2+}$ source term \\
ER release & Store-release ODE \cite{li1994equations,shannon2004mathematical} & $J_{\mathrm{ER,rel}}=k_{rel}O_{\mathrm{IP3R/RyR}}(C_i,C_{ER})$ & channel state, store content & ER-to-cytosol flux \\
Pump/exchanger recovery & Saturating fluxes \cite{shannon2004mathematical} & $J_{\mathrm{PMCA}}=V_{max}\frac{C_i^n}{K^n+C_i^n}$ & Ca$^{2+}$, Na$^+$, ATP, voltage & Ca$^{2+}$ clearance \\
Mitochondrial buffering & Coupled uptake/release ODE \cite{magnus1997minimal,bertram2006simplified} & $\frac{dC_m}{dt}=J_{\mathrm{MCU}}-J_{\mathrm{NCLX}}-J_{\mathrm{mPTP}}$ & $C_i$, $\Delta\Psi_m$, ROS & mitochondrial Ca$^{2+}$ state \\
ROS dynamics & Source-clearance ODE \cite{cortassa2004mitochondrial} & $\frac{dR}{dt}=k_{ROS}f(C_m,\Delta\Psi_m)-k_{clear}R$ & mitochondrial state, antioxidants & ROS state variable \\
ATP state & Production-consumption ODE \cite{beard2005biophysical,bertram2006simplified} & $\frac{dA}{dt}=k_{prod}f(\Delta\Psi_m)-k_{use}(J_{pump}+R_{repair})$ & mitochondria, pumps, repair & ATP availability \\
\bottomrule
\end{tabular}
\rowcolors{2}{}{}
\abbrevnote{ER, endoplasmic reticulum; GHK, Goldman-Hodgkin-Katz; PMCA, plasma membrane Ca$^{2+}$ ATPase; MCU, mitochondrial calcium uniporter; NCLX, Na$^+$/Ca$^{2+}$/Li$^+$ exchanger; ROS, reactive oxygen species; ATP, adenosine triphosphate; mPTP, mitochondrial permeability transition pore.}
\end{table*}

% ============================================================
\subsection*{Ca$^{2+}$-Dependent Remodeling and Repair}
% ============================================================

This module connects global and submembrane Ca$^{2+}$ signals to local membrane remodeling and repair. Although cytosolic Ca$^{2+}$ may be represented as a well-mixed compartment in reduced models, electro-exocytosis is likely governed by local Ca$^{2+}$ domains near pores, damaged membrane regions, budding sites, repair complexes, and fusion machinery. Reaction--diffusion theory provides a practical reason to distinguish these domains from the bulk pool: mobile and immobile buffers can confine Ca$^{2+}$ elevations over nanometer-to-submicrometer distances, with local gradients forming on submillisecond time scales near an open channel \cite{naraghi1997linearized}. Ca$^{2+}$-dependent resealing occurs on experimentally measurable time scales, with nano- and microsecond electroporation studies reporting fast and slow repair components that are sensitive to extracellular Ca$^{2+}$ \cite{bhattacharya2022calcium}. Local Ca$^{2+}$ therefore serves not only as an influx variable but also as the organizing signal for lipid asymmetry, cytoskeletal mechanics, curvature, scission, fusion competence, and resealing. The paired evidence and modeling assumptions for these repair and remodeling variables are tabulated in Tables~\ref{T_A7_bio_remodel_repair} and \ref{T_A8_comp_remodel_repair}.

For the current computational framework, this layer is implemented as a reduced Ca$^{2+}$-sensor and repair-state module. The upstream ion-transport state is converted into a submembrane Ca$^{2+}$ microdomain proxy amplified by pore activation and osmotic stress, with buffer strength and shell depth treated as tunable parameters rather than fixed geometry. This choice reflects evidence that ESCRT and repair machinery assemble at damaged membrane domains rather than responding only to bulk cytosolic Ca$^{2+}$ \cite{scheffer2014mechanism,shukla2022vitro}. Annexin recruitment is represented as a Ca$^{2+}$-dependent membrane-binding and curvature-stabilization state because annexin A4/A6/A7 can support membrane constriction, ESCRT-associated repair, and repair-linked vesicle secretion \cite{boye2017annexin,sonder2019annexin,muratori2021role,williams2023annexin}. Molecular and biophysical annexin studies then provide the curvature and lipid-composition rationale for keeping annexin effects separate from generic repair activity \cite{mularski2022interplay,hakamiZanjani2024modulation}.

The lipid-remodeling branch maps local Ca$^{2+}$ to scramblase activation, flippase suppression, and PS externalization. This branch is treated as established because nsPEF activates TMEM16F-dependent PS exposure and blebbing, while newer TIRF studies show delayed, spatially punctate PS externalization after 400-ns pulses \cite{muratori2017activation,silkunas2026peculiarities}. The Vernier studies add a second, faster computationally tractable route: nanopore-facilitated, voltage-driven PS movement can occur along electropore walls and appears experimentally as localized, anode-facing PS externalization after nanosecond pulses \cite{vernier2006nanoporeJACS,vernier2006nanoporePhysBio}. The model therefore keeps an immediate pore/voltage PS term separate from slower Ca$^{2+}$--scramblase dynamics. Chemically controlled TMEM16F vesiculation, lipid-asymmetry loss, and Ca$^{2+}$-regulated flippases provide complementary support for linking Ca$^{2+}$ and PS state to EV-like budding competence \cite{han2019chemically,wang2025loss,segawa2016human}.

The cytoskeletal branch maps local Ca$^{2+}$ to calpain activity, actin disruption, actomyosin tension, and repair-associated shedding. Recent nsPEF actin and swelling studies are especially useful here because Ca$^{2+}$-dependent actin shrinkage and swelling can be tested directly against the model's actin-disruption and osmotic-stress terms \cite{hellwich2026extracellular,pakhomov2014disassembly,rassokhin2014cellular}. Calpain and ezrin-cleavage studies support the protease component, and membrane-repair studies show that calpain can orchestrate annexin-containing microvesicle release \cite{campbell2012structure,roberts2020ca,williams2025calpains}.

Finally, lysosomal repair is represented as a Ca$^{2+}$-triggered vesicle-fusion and membrane-resealing support state. Lysosomal exocytosis, extracellular protease release, sphingolipid-dependent EV biology, lysosomal-channel regulation, and ESCRT-mediated lysosome repair support this reduced repair branch \cite{andrews2014damage,castroGomes2016plasma,verderio2018role,rosato2021two,radulovic2018escrt}. The resulting variables--PS exposure, calpain activity, annexin activity, lysosomal repair, actin state, resealing state, and repair-associated shedding--serve as interpretable coupling terms for medium/large EV budding, membrane recovery, and injury-quality gating.

\begin{table*}[t]
\footnotesize\sf
\centering
\caption{Biological evidence table for Ca$^{2+}$-dependent remodeling and repair.\label{T_A7_bio_remodel_repair}}
\rowcolors{2}{gray!10}{white}
\begin{tabular}{p{2.7cm} p{3.4cm} p{4.2cm} p{3.0cm} p{2.0cm}}
\toprule
\textbf{Submodule} & \textbf{Biological phenomenon} & \textbf{Evidence to curate} & \textbf{Relevance to EV engineering} & \textbf{Confidence} \\
\midrule
Ca$^{2+}$ microdomains & Local Ca$^{2+}$ elevation near pores and membranes & High-speed/local Ca$^{2+}$ imaging, buffer sensitivity & Localizes activation to budding, repair, and fusion sites & Supported \cite{scheffer2014mechanism,shukla2022vitro} \\
Annexin recruitment & Ca$^{2+}$-dependent membrane binding & Annexin translocation and repair studies & Links Ca$^{2+}$ to membrane repair and curvature stabilization & Supported \cite{boye2017annexin,mularski2022interplay,hakamiZanjani2024modulation,muratori2021role,sonder2019annexin} \\
Scramblase/PS exposure & Loss of lipid asymmetry & Annexin V staining, scramblase/flippase studies & Promotes MV-like budding and EV surface phenotype & Established \cite{muratori2017activation,wang2025loss,han2019chemically,silkunas2026peculiarities,segawa2016human} \\
Calpain activation & Ca$^{2+}$-dependent protease activity & Calpain activation/inhibition, cytoskeletal remodeling & Supports blebbing, scission, and membrane remodeling & Supported \cite{williams2025calpains,campbell2012structure,roberts2020ca} \\
Actomyosin remodeling & Tension, contraction, and scission & RhoA/ROCK, myosin, actin perturbation & Determines budding probability and vesicle detachment & Supported \cite{hellwich2026extracellular,pakhomov2014disassembly,rassokhin2014cellular} \\
Lysosomal repair & Ca$^{2+}$-triggered repair-vesicle fusion & LAMP1 exposure, lysosome fusion & Adds membrane and may release vesicle-like material & Supported \cite{andrews2014damage,castroGomes2016plasma,verderio2018role,rosato2021two,radulovic2018escrt} \\
\bottomrule
\end{tabular}
\rowcolors{2}{}{}
\abbrevnote{EV, extracellular vesicle; PS, phosphatidylserine; ESCRT, endosomal sorting complex required for transport.}
\end{table*}

\begin{table*}[t]
\footnotesize\sf
\centering
\caption{Computational framework table for Ca$^{2+}$-dependent remodeling and repair.\label{T_A8_comp_remodel_repair}}
\rowcolors{2}{gray!10}{white}
\begin{tabular}{p{2.6cm} p{3.0cm} p{4.6cm} p{2.8cm} p{2.8cm}}
\toprule
\textbf{Submodule} & \textbf{Model form} & \textbf{Representative relation} & \textbf{Inputs} & \textbf{Outputs} \\
\midrule
Ca$^{2+}$ microdomain & Reaction--diffusion or shell ODE \cite{naraghi1997linearized,shannon2004mathematical} & $\frac{\partial C}{\partial t}=D_{Ca}\nabla^2C+J_{Ca}-R_{buffer}-R_{effector}$ & local flux, diffusion, buffers & $C_{sub}(t)$ \\
Generic Ca$^{2+}$ sensor & Hill activation \cite{naraghi1997linearized,shannon2004mathematical,scheffer2014mechanism} & $A_x=\frac{C_{sub}^n}{K_x^n+C_{sub}^n}$ & local Ca$^{2+}$, sensitivity & activated fraction \\
PS exposure & Activation plus lipid-state ODE \cite{vernier2006nanoporeJACS,vernier2006nanoporePhysBio,muratori2017activation} & $\frac{dPS_{out}}{dt}=k_{scr}A_{scr}(1-PS_{out})-k_{flip}PS_{out}$ & Ca$^{2+}$, scramblase, flippase & PS externalization \\
Calpain/cytoskeleton & Protease activation and substrate loss \cite{campbell2012structure,roberts2020ca,williams2025calpains} & $\frac{dF_{actin}}{dt}=-k_{calp}A_{calp}F_{actin}+k_{poly}$ & Ca$^{2+}$, calpain state & cortical remodeling \\
Membrane resealing & Conductance-recovery ODE \cite{tarek2005membrane,needham1989electro} & $\frac{dG_m}{dt}=-k_{seal}A_{repair}G_m$ & residual conductance, repair activity & resealed membrane state \\
Damaged-domain shedding & Damage-removal rate \cite{needham1989electro} & $R_{shed}=k_{shed}D_m f_{Ca}f_{tension}$ & damaged membrane, Ca$^{2+}$, tension & repair-associated shedding \\
\bottomrule
\end{tabular}
\rowcolors{2}{}{}
\abbrevnote{PS, phosphatidylserine; ODE, ordinary differential equation.}
\end{table*}

% ============================================================
\subsection*{EV Biogenesis and Subtype Release}
% ============================================================

The EV-release module converts nsPEF-induced intracellular and membrane responses into measurable vesicle outputs. It integrates endosomal/MVB-derived small-EV release, plasma-membrane budding of medium/large EVs, and apoptotic body or injury-associated vesiculation. The model should retain mechanistic terms such as exosome and microvesicle internally, but report outputs using EV nomenclature compatible with the level of experimental characterization. The biological release mechanisms are organized in Table~\ref{T_A9_bio_ev_release}, with candidate pool-balance and rate-law implementations listed in Table~\ref{T_A10_comp_ev_release}.

Small-EV production begins with endosome maturation and ILV biogenesis. Recent Rab-conversion studies indicate that Rab5-to-Rab7 transition is not an isolated timing event, but is coordinated with ESCRT-mediated cargo corralling through the Mon1/Ccz1 axis, so maturation and secretory competence should be treated as coupled states rather than as independent switches \cite{ott2025coordination,borchers2023regulatory,solinger2025escrting}. Within maturing endosomes, canonical ESCRT machinery supplies one major ILV route: ESCRT-0 captures ubiquitinated cargo through Hrs/STAM, ESCRT-I/II concentrate cargo and deform the bud neck, and ESCRT-III/Vps4 completes constriction and scission \cite{takahashi2015hrs,wollert2010molecular,pfitzner2020escrt}. However, exosome-biased ILV formation is not exclusively canonical. ALIX-ESCRT-III, ceramide, and tetraspanin-rich membrane domains provide selective alternative routes that can redirect cargo and alter whether MVBs are routed toward secretion or lysosomal degradation \cite{larios2020alix,choezom2022nsmase2,crivelli2022ceramideTransfer,rubinstein2025tetraspanins,villarroyaBeltri2016isgylation}.

Release of those ILVs as extracellular small EVs then depends on trafficking, docking, and fusion at the cell cortex. Rab27, Rab11, and Rab35 families regulate MVB positioning and membrane access, while Munc13-4 provides a Ca$^{2+}$-responsive trafficking and priming branch that helps generate release-competent MVBs \cite{hsu2010rab35,messenger2018munc13,marie2023accessory,izumi2021rab27,tang2026rab27}. The terminal fusion step is mediated by cell-context-dependent SNARE assemblies, with syntaxin-4/SNAP-23/VAMP-7 and VAMP5-based complexes now supported by direct exosome-secretion studies, and with synaptotagmin-7-like Ca$^{2+}$ sensing providing a plausible fusion-pore stabilization mechanism under strong submembrane Ca$^{2+}$ signals \cite{liu2023snare,matsui2023vamp5,woo2017munc13,courtney2023synaptotagmin7}. In parallel, outward budding of medium/large EVs is supported by ARF-, RhoA/ROCK-, phosphatidylserine-, and actomyosin-dependent membrane remodeling, whereas apoptotic bodies arise from caspase-linked ROCK1 activation and repetitive bleb fragmentation during higher-injury states \cite{pang2025arf,dai2019rhoa,menck2017nsmaseMicrovesicles,teixeira2020rock1,ozdemir2021rock1,sakamaki2025mprip}. These mechanistic branches justify treating subtype release as a competition between secretory MVB fate, plasma-membrane budding competence, and apoptotic disassembly rather than as one generic EV-rate variable.

\begin{table*}[t]
\footnotesize\sf
\centering
\caption{Biological evidence table for EV biogenesis and subtype release.\label{T_A9_bio_ev_release}}
\rowcolors{2}{gray!10}{white}
\begin{tabular}{p{2.7cm} p{3.4cm} p{4.2cm} p{3.0cm} p{2.0cm}}
\toprule
\textbf{Submodule} & \textbf{Biological phenomenon} & \textbf{Evidence to curate} & \textbf{Relevance to EV engineering} & \textbf{Confidence} \\
\midrule
Endosome/MVB maturation & Early-to-late endosome and MVB formation & Rab5/Rab7, endosome maturation, ILV evidence & Defines small-EV-producing compartment pool & Established \cite{ott2025coordination,borchers2023regulatory,solinger2025escrting} \\
ESCRT-dependent ILV formation & Cargo sorting and ILV biogenesis & ESCRT, ALIX, TSG101 perturbation & Generates exosome-enriched precursors & Established \cite{takahashi2015hrs,wollert2010molecular,pfitzner2020escrt,larios2020alix} \\
ESCRT-independent ILV formation & Ceramide/tetraspanin/lipid-domain sorting & nSMase2, ceramide, CD63/CD81/CD9 evidence & Alternative small-EV biogenesis route & Supported \cite{choezom2022nsmase2,crivelli2022ceramideTransfer,rubinstein2025tetraspanins} \\
MVB trafficking/docking & Rab-dependent transport to plasma membrane & Rab27A/B, Rab11, Rab35 evidence & Controls release-ready MVB pool & Established \cite{hsu2010rab35,messenger2018munc13,izumi2021rab27,tang2026rab27,marie2023accessory} \\
MVB fusion & SNARE-mediated fusion with plasma membrane & SNARE and Ca$^{2+}$-sensitive fusion evidence & Converts ILVs into secreted small EVs & Supported \cite{liu2023snare,matsui2023vamp5,woo2017munc13,courtney2023synaptotagmin7} \\
Plasma membrane budding & MV-like outward budding & ARF6/RhoA/PS/cytoskeleton studies & Produces medium/large EV precursors & Supported \cite{pang2025arf,dai2019rhoa,menck2017nsmaseMicrovesicles} \\
Apoptotic body formation & Blebbing and fragmentation & Caspase/ROCK1/apoptosis evidence & Captures high-dose or injury-associated vesicles & Established \cite{teixeira2020rock1,ozdemir2021rock1,sakamaki2025mprip} \\
\bottomrule
\end{tabular}
\rowcolors{2}{}{}
\abbrevnote{EV, extracellular vesicle; MVB, multivesicular body; ILV, intraluminal vesicle; ESCRT, endosomal sorting complex required for transport.}
\end{table*}

\begin{table*}[t]
\footnotesize\sf
\centering
\caption{Computational framework table for EV biogenesis and subtype release.\label{T_A10_comp_ev_release}}
\rowcolors{2}{gray!10}{white}
\begin{tabular}{p{2.6cm} p{3.0cm} p{4.6cm} p{2.8cm} p{2.8cm}}
\toprule
\textbf{Submodule} & \textbf{Model form} & \textbf{Representative relation} & \textbf{Inputs} & \textbf{Outputs} \\
\midrule
MVB pool & Compartmental pool balance \cite{huotari2011endosome,ott2025coordination,borchers2023regulatory,solinger2025escrting} & $\frac{dM}{dt}=k_{mat}E_{endo}-k_{dock}M-k_{deg}M$ & endosome pool, maturation & MVB number \\
ILV loading & Production-rate law \cite{takahashi2015hrs,wollert2010molecular,pfitzner2020escrt,larios2020alix,choezom2022nsmase2,crivelli2022ceramideTransfer} & $\frac{dI}{dt}=k_{ILV}f_{ESCRT}M-k_{release}I$ & ESCRT/ceramide state, MVB pool & ILV load \\
MVB docking & Trafficking/docking ODE \cite{hsu2010rab35,messenger2018munc13,izumi2021rab27,tang2026rab27,marie2023accessory} & $\frac{dM_d}{dt}=k_{dock}f_{Rab}M-k_{fus}M_d$ & Rab activity, actin barrier & docked MVB pool \\
Small-EV release & Ca$^{2+}$-sensitive fusion rate \cite{liu2023snare,matsui2023vamp5,woo2017munc13,courtney2023synaptotagmin7} & $R_{sEV}=k_{fus}f_{Ca}(C_{sub})M_d$ & docked MVBs, Ca$^{2+}$, SNARE & small-EV release rate \\
Budding-site formation & Budding pool ODE \cite{pang2025arf,dai2019rhoa,menck2017nsmaseMicrovesicles} & $\frac{dB_{m/lEV}}{dt}=k_{bud}f_{PS}f_{tension}-k_{scission}B_{m/lEV}$ & PS, curvature, tension & budding precursor pool \\
Medium/large EV release & Scission-rate law \cite{pang2025arf,dai2019rhoa,menck2017nsmaseMicrovesicles,williams2025calpains} & $R_{m/lEV}=k_{scission}f_{actomyosin}B_{m/lEV}$ & calpain, ROCK, actomyosin & medium/large EV rate \\
Apoptotic vesicle release & Injury-state weighted rate \cite{teixeira2020rock1,ozdemir2021rock1,sakamaki2025mprip,cappe2023systematic} & $R_{AB}=k_{AB}F_{apop}$ & apoptotic fraction, blebbing rate & apoptotic-body release \\
\bottomrule
\end{tabular}
\rowcolors{2}{}{}
\abbrevnote{MVB, multivesicular body; ILV, intraluminal vesicle; ESCRT, endosomal sorting complex required for transport; sEV, small extracellular vesicle; m/lEV, medium/large extracellular vesicle; AB, apoptotic body.}
\end{table*}

For the current computational framework, this layer is implemented as a reduced pool-balance model with explicit MVB, ILV, docked-MVB, budding-precursor, and apoptotic-commitment states. Endosome maturation is represented as a Rab-conversion-supported source term for a secretory MVB pool, ILV loading is split into ESCRT-dependent and ceramide/tetraspanin-assisted branches, and MVB fate is biased between secretion and lysosomal loss by an acidification-sensitive routing term \cite{ott2025coordination,choezom2022nsmase2,villarroyaBeltri2016isgylation}. Small-EV release is then written as a Ca$^{2+}$- and SNARE-sensitive fusion flux from the docked MVB pool, while medium/large EV release is represented as a phosphatidylserine-, calpain-, and actomyosin-weighted budding/scission process. Apoptotic-body release is modeled separately as a damage- and blebbing-weighted branch so that injury-associated vesiculation is not conflated with productive secretory EV release \cite{liu2023snare,messenger2018munc13,pang2025arf,teixeira2020rock1}. This reduced formulation matches the current code implementation while preserving a direct path for later replacement with cell-type-specific Rab, SNARE, lipid, or apoptosis submodels as quantitative data become available.

\input{content/size_resolved_bridge}

% ============================================================
\subsection*{Cargo Sorting, EV Composition, and Potency}
% ============================================================

The cargo module represents the major conceptual expansion from EV-release modeling to EV-engineering modeling. nsPEF may alter EV composition through Ca$^{2+}$-dependent trafficking, stress-responsive transcription and translation, lipid remodeling, ESCRT/tetraspanin/ceramide pathways, RNA-binding proteins, and injury-state contamination. Therefore, EV output should be represented as a multi-dimensional product state rather than as total particle number alone. Table~\ref{T_A11_bio_cargo} separates the biological cargo-sorting evidence, and Table~\ref{T_A12_comp_cargo} specifies the computational cargo, composition, potency, and direct-loading model options.

\begin{table*}[t]
\footnotesize\sf
\centering
\caption{Biological evidence table for cargo sorting, EV composition, and potency.\label{T_A11_bio_cargo}}
\rowcolors{2}{gray!10}{white}
\begin{tabular}{p{2.7cm} p{3.4cm} p{4.2cm} p{3.0cm} p{2.0cm}}
\toprule
\textbf{Submodule} & \textbf{Biological phenomenon} & \textbf{Evidence to curate} & \textbf{Relevance to EV engineering} & \textbf{Confidence} \\
\midrule
Protein sorting & Ubiquitin/ESCRT/tetraspanin-dependent protein enrichment & ALIX, TSG101, CD63/CD81/CD9, antigen enrichment & Determines therapeutic or immunogenic protein cargo & Established \cite{jeppesen2019reassessment,kugeratski2021quantitative,takahashi2015hrs,larios2020alix,rubinstein2025tetraspanins} \\
RNA sorting & RNA-binding protein-mediated miRNA/mRNA loading & hnRNPA2B1, SYNCRIP, YBX1, motif-based sorting & Enables nucleic-acid cargo modulation & Supported \cite{shurtleff2016ybox,liu2021selective,santangelo2016syncrip} \\
Lipid sorting & Ceramide, cholesterol, sphingolipid, PE/PS enrichment & Lipidomics, nSMase2, tetraspanin domains & Controls EV stability, curvature, uptake, immunologic function & Established \cite{skotland2023lipids,choezom2022nsmase2,crivelli2022ceramideTransfer,rubinstein2025tetraspanins} \\
Stress-responsive cargo & ROS, hypoxia, inflammatory, apoptotic cargo changes & proteomics/lipidomics/RNA-seq after nsPEF & Distinguishes programmed cargo modulation from debris & Supported \cite{pakhomova2012oxidative,cappe2023systematic,muratori2017activation} \\
Antigen enrichment & Tumor antigen expression and EV loading & MAGE and other antigen studies & Supports vaccine-oriented EV engineering & Supported \cite{szlasa2024mage,kumar2022tumour} \\
Functional potency & Recipient-cell response & uptake, cargo delivery, immune activation, repair assays & Connects cargo state to biological activity & Hypothesized \cite{hagey2023cellular,tkach2017qualitative,fuhrmann2015active} \\
\bottomrule
\end{tabular}
\rowcolors{2}{}{}
\abbrevnote{EV, extracellular vesicle; ESCRT, endosomal sorting complex required for transport; ALIX, ALG-2-interacting protein X; TSG101, tumor susceptibility gene 101; miRNA, microRNA; mRNA, messenger RNA; PE, phosphatidylethanolamine; PS, phosphatidylserine; RNA-seq, RNA sequencing.}
\end{table*}

\begin{table*}[t]
\footnotesize\sf
\centering
\caption{Computational framework table for cargo sorting, EV composition, and potency.\label{T_A12_comp_cargo}}
\rowcolors{2}{gray!10}{white}
\begin{tabular}{p{2.6cm} p{3.0cm} p{4.6cm} p{2.8cm} p{2.8cm}}
\toprule
\textbf{Submodule} & \textbf{Model form} & \textbf{Representative relation} & \textbf{Inputs} & \textbf{Outputs} \\
\midrule
Subtype-weighted cargo & Weighted output model \cite{jeppesen2019reassessment,kugeratski2021quantitative,cappe2023systematic} & $Cargo_j=w_{j,sEV}R_{sEV}+w_{j,m/lEV}R_{m/lEV}+w_{j,AB}R_{AB}$ & subtype rates, cargo weights & cargo abundance \\
Stress-modulated sorting & Rate law \cite{shurtleff2016ybox,liu2021selective,santangelo2016syncrip,pakhomova2012oxidative} & $\frac{dCargo_j}{dt}=k_{sort,j}f_{stress}f_{Ca}f_{ESCRT}-k_{loss,j}Cargo_j$ & Ca$^{2+}$, ROS, ESCRT, stress & cargo state \\
Antigen enrichment & Transcription/sorting coupling \cite{szlasa2024mage,kumar2022tumour} & $\frac{dA_g}{dt}=k_{expr}f_{stress}-k_{deg}A_g-k_{sort}A_g$ & stress, antigen expression & antigen cargo \\
Lipid composition & Composition-state ODE \cite{skotland2023lipids,choezom2022nsmase2,crivelli2022ceramideTransfer} & $\frac{dL_j}{dt}=k_{syn,j}f_{stress}-k_{sort,j}L_j$ & lipid metabolism, EV subtype & lipid cargo class \\
Potency mapping & Empirical or mechanistic response model \cite{hagey2023cellular,tkach2017qualitative,fuhrmann2015active} & $Potency=F(Cargo,Subtype,Dose)$ & cargo vector, recipient assay & functional potency \\
Direct EV loading & Permeability-diffusion model \cite{kooijmans2013electroporation,hood2014maximizing,fuhrmann2015active,didiot2018optimized} & $\frac{dC_{load}}{dt}=P_{EV}(t)A_{EV}(C_{ext}-C_{EV})-k_{leak}C_{EV}$ & EV permeability, cargo gradient & loaded cargo \\
\bottomrule
\end{tabular}
\rowcolors{2}{}{}
\abbrevnote{EV, extracellular vesicle; ESCRT, endosomal sorting complex required for transport; ALIX, ALG-2-interacting protein X; sEV, small extracellular vesicle; m/lEV, medium/large extracellular vesicle; AB, apoptotic body.}
\end{table*}


The implemented Layer 6 module now follows the table structure explicitly. It computes a subtype-weighted cargo state from small-EV, medium/large-EV, and apoptotic-body outputs; separates protein, RNA, lipid, antigen, and direct-loading terms; and maps the resulting cargo vector to a saturating potency score. This keeps the implementation aligned with evidence that EV subtypes carry distinct protein compositions \cite{jeppesen2019reassessment,kugeratski2021quantitative}, that RNA sorting can be mediated by RNA-binding proteins such as YBX1 and SYNCRIP \cite{shurtleff2016ybox,liu2021selective,santangelo2016syncrip}, that lipid composition and ceramide biology shape EV state \cite{skotland2023lipids}, and that direct EV loading can be represented separately from endogenous cargo sorting \cite{fuhrmann2015active,didiot2018optimized}.

% ============================================================
\subsection*{Injury, Debris, and Quality Gate}
% ============================================================

The injury-quality module defines whether nsPEF exposure produces reversible EV engineering, stress-associated release, apoptosis, necrosis, or contaminating debris. This layer is essential because increased particle number is not necessarily equivalent to increased useful EV production. The quality gate separates bona fide EV output from apoptotic bodies, damaged organellar fragments, protein aggregates, and other co-isolated materials. Table~\ref{T_A13_bio_quality_gate} summarizes the biological and QC evidence for this distinction, while Table~\ref{T_A14_comp_quality_gate} defines the damage, fate-transition, mixture, and constrained-optimization model forms.

\begin{table*}[t]
\footnotesize\sf
\centering
\caption{Biological evidence table for injury, debris, and quality gating.\label{T_A13_bio_quality_gate}}
\rowcolors{2}{gray!10}{white}
\begin{tabular}{p{2.7cm} p{3.4cm} p{4.2cm} p{3.0cm} p{2.0cm}}
\toprule
\textbf{Submodule} & \textbf{Biological phenomenon} & \textbf{Evidence to curate} & \textbf{Relevance to EV engineering} & \textbf{Confidence} \\
\midrule
Viability preservation & Reversible response versus cell death & viability, membrane recovery, clonogenic/metabolic assays & Defines acceptable EV-production window & Established \cite{pakhomova2013two,beebe2012transient,bhattacharya2022calcium} \\
Mitochondrial injury & Depolarization, mPTP, cytochrome c & $\Delta\Psi_m$, cytochrome c, ROS & Separates stress signaling from organellar contamination & Supported \cite{beebe2012transient,napotnik2012mitochondrial,pakhomova2012oxidative} \\
Apoptosis & Caspase activation and apoptotic blebbing & caspase, annexin V/PI, ROCK1 blebbing & Determines apoptotic-body contribution & Established \cite{pakhomova2013two,teixeira2020rock1,cappe2023systematic} \\
Necrosis/lysis & Irreversible membrane failure & LDH release, PI uptake, debris microscopy & Indicates unacceptable product contamination & Established \cite{pakhomova2013two,cappe2023systematic} \\
Non-EV contaminants & Protein aggregates, organelle fragments, apoptotic bodies & calnexin, cytochrome c, albumin/lipoprotein markers & Defines purity and interpretability of EV readouts & Established \cite{webber2013pure,brennan2020comparison,cappe2023systematic,welsh2024minimal} \\
Quality-control marker panel & Positive and negative EV markers & CD63/CD81/CD9, TSG101/ALIX/syntenin, negative markers & Connects model to MISEV-style QC & Established \cite{welsh2024minimal,jeppesen2019reassessment,kugeratski2021quantitative} \\
\bottomrule
\end{tabular}
\rowcolors{2}{}{}
\abbrevnote{EV, extracellular vesicle; mPTP, mitochondrial permeability transition pore; PI, propidium iodide; QC, quality control; MISEV, Minimal Information for Studies of Extracellular Vesicles.}
\end{table*}

\begin{table*}[t]
\footnotesize\sf
\centering
\caption{Computational framework table for injury, debris, and quality gating.\label{T_A14_comp_quality_gate}}
\rowcolors{2}{gray!10}{white}
\begin{tabular}{p{2.6cm} p{3.0cm} p{4.6cm} p{2.8cm} p{2.8cm}}
\toprule
\textbf{Submodule} & \textbf{Model form} & \textbf{Representative relation} & \textbf{Inputs} & \textbf{Outputs} \\
\midrule
Damage accumulation & Damage-state ODE \cite{pakhomova2012oxidative,bhattacharya2022calcium,albeck2008modeling} & $\frac{dD}{dt}=k_{inj}G_m+k_{ROS}ROS-k_{repair}A_{repair}D$ & conductance, ROS, repair & damage burden \\
Fate transition & Hybrid state rule \cite{albeck2008modeling,bentele2004mathematical,pakhomova2013two} & $F_{apop}=H(D-D^*)$ or $\frac{dF_{apop}}{dt}=k_{apop}H(D-D^*)$ & damage, ATP, mitochondria & viable/stressed/apoptotic fractions \\
Necrotic transition & Threshold rule \cite{pakhomova2013two,cappe2023systematic} & $F_{nec}=H(D-D_{nec}^*)H(A_{ATP}^*-ATP)$ & damage, ATP, membrane failure & necrotic fraction \\
Measured particles & Mixture model \cite{cappe2023systematic,webber2013pure,brennan2020comparison} & $N_{meas}=N_{EV}+N_{AB}+N_{debris}+N_{agg}$ & release rates, injury state & measured particle count \\
Purity score & Ratio or classifier \cite{webber2013pure,brennan2020comparison,welsh2024minimal} & $Q_{EV}=\frac{N_{bona\ fide\ EV}}{N_{meas}}$ & marker panel, debris state & EV quality score \\
Acceptable window & Constrained optimization \cite{orlacchio2023effects,gudvangen2022electroporation,welsh2024minimal} & maximize $Y_{EV}$ subject to $Q_{EV}>Q^*$ and $Viability>V^*$ & dose, cell state, harvest & operating region \\
\bottomrule
\end{tabular}
\rowcolors{2}{}{}
\abbrevnote{EV, extracellular vesicle; AB, apoptotic body; ATP, adenosine triphosphate; ROS, reactive oxygen species; QC, quality control.}
\end{table*}


The implemented Layer 7 module expands the previous binary quality proxy into a mixture-aware gate. Damage is mapped to stressed-viable, apoptotic, and necrotic fractions using threshold-like Hill terms inspired by systems models of cell-death switching \cite{albeck2008modeling,bentele2004mathematical}. The measured particle pool is then decomposed into bona fide EVs, apoptotic bodies, debris, and aggregate-like contaminants, which makes the purity score sensitive to both injury state and non-EV co-isolates. This structure matches EV-composition and isolation cautions from apoptotic/necroptotic EV studies and MISEV-style purity guidance \cite{cappe2023systematic,webber2013pure,brennan2020comparison,welsh2024minimal}.

% ============================================================
\subsection*{Manufacturing, Isolation, and Quality-Control Interface}
% ============================================================

The manufacturing module links cell-level mechanisms to translational EV production. It represents harvest timing, cell number, culture format, isolation recovery, marker-based QC, purity, batch variability, and flow-through scalability. This module is intentionally placed after biological release and quality gating because manufacturing success depends on both biological production and downstream process recovery. Table~\ref{T_A15_bio_manufacturing} summarizes the process and QC evidence to curate, and Table~\ref{T_A16_comp_manufacturing} lists the process-yield, recovery, batch, and optimization model forms.

\begin{table*}[t]
\footnotesize\sf
\centering
\caption{Biological/process evidence table for manufacturing, isolation, and quality-control interface.\label{T_A15_bio_manufacturing}}
\rowcolors{2}{gray!10}{white}
\begin{tabular}{p{2.7cm} p{3.4cm} p{4.2cm} p{3.0cm} p{2.0cm}}
\toprule
\textbf{Submodule} & \textbf{Process phenomenon} & \textbf{Evidence to curate} & \textbf{Relevance to EV engineering} & \textbf{Confidence} \\
\midrule
Harvest kinetics & Time-dependent EV accumulation & minute-to-hour release curves after nsPEF & Determines optimal collection window & Supported \cite{orlacchio2023effects,gudvangen2022electroporation,mendt2018generation} \\
Cell-normalized yield & EV output per cell or protein & NTA/flow/ELISA normalized to cell number & Defines productivity & Established \cite{welsh2024minimal,paganini2019scalable,mendt2018generation} \\
Isolation recovery & SEC, ultracentrifugation, density gradient, affinity capture & recovery and purity comparisons & Converts released EVs into isolated product & Established \cite{watson2018scalable,brennan2020comparison,webber2013pure} \\
QC characterization & Size, morphology, markers, contaminants & NTA, cryo-EM/TEM, western blot/flow, proteomics & Determines whether product meets EV criteria & Established \cite{welsh2024minimal,jeppesen2019reassessment,kugeratski2021quantitative} \\
Batch variability & Donor, passage, culture, and pulse variability & repeated-batch studies, mixed-effects data & Determines reproducibility and manufacturability & Supported \cite{paganini2019scalable,watson2018scalable,mendt2018generation,thakur2024global} \\
Scalable flow processing & Closed-system pulse chamber and inline monitoring & field uniformity, sterility, temperature, flow control & Enables translational manufacturing & Hypothesized \cite{liu2026microfluidic,garcia2016microfluidic,bonakdar2017microfluidic,thakur2024global} \\
\bottomrule
\end{tabular}
\rowcolors{2}{}{}
\abbrevnote{EV, extracellular vesicle; nsPEF, nanosecond pulsed electric field; NTA, nanoparticle tracking analysis; ELISA, enzyme-linked immunosorbent assay; QC, quality control; TEM, transmission electron microscopy.}
\end{table*}


The implemented Layer 8 module converts biological output into process-facing quantities: cell-normalized yield, harvest-rate-equivalent accumulation, isolation recovery, process purity, batch-adjusted yield, and a constrained optimization objective. This keeps the code aligned with EV manufacturing literature emphasizing scalable source material, recovery and purity tradeoffs, clinical-grade purification, and validation requirements \cite{paganini2019scalable,watson2018scalable,mendt2018generation,thakur2024global,welsh2024minimal}. The objective is intentionally transparent rather than prescriptive: potency, yield, purity, and viability weights can be changed for different process-development priorities.

\begin{table*}[t]
\footnotesize\sf
\centering
\caption{Computational framework table for manufacturing, isolation, and quality-control interface.\label{T_A16_comp_manufacturing}}
\rowcolors{2}{gray!10}{white}
\begin{tabular}{p{2.6cm} p{3.0cm} p{4.6cm} p{2.8cm} p{2.8cm}}
\toprule
\textbf{Submodule} & \textbf{Model form} & \textbf{Representative relation} & \textbf{Inputs} & \textbf{Outputs} \\
\midrule
Cell-normalized yield & Productivity metric \cite{welsh2024minimal,paganini2019scalable,mendt2018generation} & $Y_{EV}=\frac{N_{EV}}{N_{cells}}$ & EV count, cell number & EV yield per cell \\
Harvest accumulation & Time-integral model \cite{orlacchio2023effects,gudvangen2022electroporation,mendt2018generation} & $N_{EV}(t_h)=\int_0^{t_h}R_{EV}(t)dt$ & release rate, harvest time & collected EVs \\
Isolation recovery & Process-yield balance \cite{watson2018scalable,brennan2020comparison,webber2013pure} & $N_{iso}=\eta_{iso}N_{released}$ & isolation method, recovery & isolated EV count \\
Purity metric & Product-quality ratio \cite{webber2013pure,brennan2020comparison,welsh2024minimal} & $Purity=\frac{N_{EV}}{Protein_{total}}$ or marker-based score & particle/protein/marker data & purity estimate \\
Batch model & Mixed-effects model \cite{paganini2019scalable,watson2018scalable,mendt2018generation,thakur2024global} & $Y_{b}=\mu+\alpha_{cell}+\beta_{dose}+u_b+\epsilon$ & cell state, dose, batch & reproducibility \\
Optimization objective & Constrained process optimization \cite{paganini2019scalable,watson2018scalable,mendt2018generation,thakur2024global,welsh2024minimal} & maximize $Potency\times Y_{iso}$ subject to purity and viability constraints & dose, harvest, isolation & optimal protocol \\
\bottomrule
\end{tabular}
\rowcolors{2}{}{}
\abbrevnote{EV, extracellular vesicle; QC, quality control.}
\end{table*}

% ============================================================
\subsection*{Cell-State and Disease-State Modifiers as Cross-Cutting Parameters}
% ============================================================

Cell-state and disease-state modifiers alter how a given cell responds to the same nsPEF exposure. In this framework, they do not define a separate downstream pathway. Instead, they scale or reshape the behavior of all modules. Relevant modifiers include membrane composition, cholesterol and sphingolipid content, cortical stiffness, cytoskeletal prestress, baseline Ca$^{2+}$ handling, mitochondrial state, metabolic capacity, oxidative-stress tone, inflammatory signaling, hypoxia response, baseline EV-biogenesis activity, and cargo-sorting phenotype. Because these modifiers act across layers, their evidence and parameter-scaling implications are distributed across the paired module tables summarized in Table~\ref{T_electro_exo_EV} and specified in Tables~\ref{T_A1_bio_dosimetry}--\ref{T_A16_comp_manufacturing}.

\begin{equation}
\theta^{cell}=S_{cell}\theta^{base},
\end{equation}

\begin{equation}
\theta_i\sim \mathcal{D}(\mu_{cell},\Sigma_{cell}),
\end{equation}

where $\theta^{cell}$ represents cell-specific parameter values, $S_{cell}$ represents phenotype-specific scaling factors, and $\mathcal{D}(\mu_{cell},\Sigma_{cell})$ represents population heterogeneity.

% ============================================================
\subsection*{Core Governing Equations for a First Integrated Implementation}
% ============================================================

A first reduced implementation can be constructed by coupling pulse exposure, membrane permeability, Ca$^{2+}$ balance, remodeling activation, EV-release rates, cargo state, repair, and quality gating. The following equations are placeholders for model assembly and should be refined as literature-supported flux expressions and parameters are identified. They assemble the computational forms listed in Tables~\ref{T_A2_comp_dosimetry}, \ref{T_A4_comp_electrodynamics}, \ref{T_A6_comp_ca_ros}, \ref{T_A8_comp_remodel_repair}, \ref{T_A10_comp_ev_release}, \ref{T_A12_comp_cargo}, and \ref{T_A14_comp_quality_gate}; manufacturing-facing optimization is summarized separately in Table~\ref{T_A16_comp_manufacturing}.

\begin{equation}
C_m\frac{d\Delta V_m}{dt}+G_m(t)\Delta V_m=f(E(t),\theta),
\end{equation}

\begin{equation}
C_o\frac{d\Delta V_o}{dt}+G_o(t)\Delta V_o=f_o(E(t),r_o,\theta_o),
\end{equation}

\begin{equation}
\frac{dN_p}{dt}=k_{form}(\Delta V)-k_{reseal}N_p,
\end{equation}

\begin{equation}
J_{i,\mathrm{pore}}=P_i(t)\Phi_i([i]_e,[i]_i,V_m),
\end{equation}

\begin{equation}
\begin{aligned}
\frac{d[\mathrm{Ca}^{2+}]_i}{dt} &=
J_{\mathrm{Ca,pore}}+J_{\mathrm{VGCC}}+J_{\mathrm{ER,rel}}
-J_{\mathrm{SERCA}} \\
&\quad -J_{\mathrm{PMCA}}-J_{\mathrm{NCX}}-J_{\mathrm{mito}}-J_{\mathrm{buffer}},
\end{aligned}
\end{equation}

\begin{equation}
A_x=\frac{[\mathrm{Ca}^{2+}]_{sub}^n}{K_x^n+[\mathrm{Ca}^{2+}]_{sub}^n},
\end{equation}

\begin{equation}
R_{EV}=R_{sEV}+R_{m/lEV}+R_{AB},
\end{equation}

\begin{equation}
R_{sEV}=k_{fus}f_{Ca}([\mathrm{Ca}^{2+}]_{sub})M_d,
\end{equation}

\begin{equation}
R_{m/lEV}=k_{scission}f_{PS}(PS_{out})f_{actomyosin}B_{m/lEV},
\end{equation}

\begin{equation}
Cargo_j=w_{j,sEV}R_{sEV}+w_{j,m/lEV}R_{m/lEV}+w_{j,AB}R_{AB},
\end{equation}

\begin{equation}
\frac{dD}{dt}=k_{inj}G_m+k_{ROS}ROS-k_{repair}A_{repair}D,
\end{equation}

\begin{equation}
N_{meas}=N_{EV}+N_{AB}+N_{debris}+N_{agg},
\end{equation}

\begin{equation}
\mathcal{Y}_{EV}=\{N_{EV},R_{EV},Subtype,Cargo,Potency,Purity,Viability\}.
\end{equation}

% ============================================================
\subsection*{Variable and Parameter Glossary Template}
% ============================================================

Table~\ref{T_A17_glossary} provides a shared notation layer for the variables and parameters introduced across the appendix tables and governing-equation scaffold.

\begin{table*}[t]
\footnotesize\sf
\centering
\caption{Variable and parameter glossary for the nsPEF-EV engineering model.\label{T_A17_glossary}}
\rowcolors{2}{gray!10}{white}
\begin{tabular}{p{2.5cm} p{5.0cm} p{2.1cm} p{2.7cm} p{3.2cm}}
\toprule
\textbf{Symbol} & \textbf{Definition} & \textbf{Units} & \textbf{Module} & \textbf{Source/estimation strategy} \\
\midrule
$E(t)$ & Applied electric-field waveform & V/m & Dosimetry & Pulse protocol/field model \\
$D_E$ & Electrical dose index & variable & Dosimetry & Derived from pulse protocol \\
$\Delta V_m$ & Plasma-membrane transmembrane potential & V & Electrodynamics & Circuit/electroporation model \\
$\Delta V_o$ & Organelle membrane potential perturbation & V & Organelle electrodynamics & Multi-shell/circuit model \\
$N_p$ & Pore density or pore number & pores/$\mu$m$^2$ or count & Electroporation & Fit or literature prior \\
$r_p$ & Representative pore radius & nm & Electroporation & Literature prior or fit \\
$G_m$ & Effective membrane conductance & S/m$^2$ & Membrane/repair & Fit from permeability recovery \\
$P_i$ & Ion-specific membrane permeability & m/s & Ion transport & Estimated from flux data \\
$J_{\mathrm{Ca,pore}}$ & Electropore-mediated Ca$^{2+}$ flux & concentration/time & Ion transport & Derived from permeability relation \\
$[\mathrm{Ca}^{2+}]_i$ & Cytosolic Ca$^{2+}$ concentration & M & Ca$^{2+}$ signaling & Ca$^{2+}$ imaging/model state \\
$[\mathrm{Ca}^{2+}]_{sub}$ & Submembrane Ca$^{2+}$ concentration & M & Microdomain & Reaction--diffusion/shell model \\
$ROS$ & Reactive oxygen species state & variable & Stress/bioenergetics & ROS imaging/probes \\
$ATP$ & Energy availability state & variable & Bioenergetics & ATP assays/model state \\
$PS_{out}$ & Externalized phosphatidylserine fraction & dimensionless & Remodeling & Annexin V or model state \\
$M_d$ & Docked MVB pool & count or concentration & EV release & Fit/estimated latent state \\
$B_{m/lEV}$ & Medium/large EV budding precursor pool & count or concentration & EV release & Fit/estimated latent state \\
$R_{sEV}$ & Small-EV release rate & particles/time & EV output & NTA/flow/marker-based fit \\
$R_{m/lEV}$ & Medium/large EV release rate & particles/time & EV output & NTA/flow/marker-based fit \\
$R_{AB}$ & Apoptotic body release rate & particles/time & Injury output & Apoptosis-weighted output \\
$Cargo_j$ & Abundance of cargo class $j$ & variable & Cargo & Proteomics/RNA/lipidomics \\
$Potency$ & Functional activity metric & assay-specific & Potency & Recipient-cell assay \\
$D$ & Integrated damage burden & dimensionless & Quality gate & Hybrid state variable \\
$Q_{EV}$ & EV quality/purity score & dimensionless & QC & Marker/purity model \\
$\eta_{iso}$ & Isolation recovery efficiency & dimensionless & Manufacturing & Process validation \\
\bottomrule
\end{tabular}
\rowcolors{2}{}{}
\abbrevnote{nsPEF, nanosecond pulsed electric field; EV, extracellular vesicle; ROS, reactive oxygen species; ATP, adenosine triphosphate; PS, phosphatidylserine; MVB, multivesicular body; ILV, intraluminal vesicle; QC, quality control.}
\end{table*}

\bibliographystyle{plain}
\bibliography{source}

\end{document}
